\documentclass[UTF8, 12pt, fontset=fandol, oneside]{ctexart}
\usepackage{researchnotes}

% 保留分册独立编译时的章号前缀。
\renewcommand{\thesection}{3.\arabic{section}}

\title{第三章\quad 函数的极限与连续性}
\author{}
\date{}

\begin{document}

\maketitle

\section{函数的极限}

上一章我们讨论了序列的极限，本节我们将进一步论述函数的极限. 若将序列看成是定义在 $\mathbb{N}$ 中的函数，则在序列极限中，自变量只有一种变化状态. 而在函数极限中的自变量却有六种不同的变化状态，这就导致函数极限的定义有各种不同的形式. 为了避免繁琐，本节我们主要以自变量趋向于某一给定的点为规范来论述函数极限的定义和性质。

\subsection{函数极限的定义}

在很多实际应用中，我们常常要分析函数的变化情况. 特别地，当函数 $f(x)$ 在某点 $x_0$ 没有意义，但在 $U_0(x_0,\delta)$ 内有意义时，我们要研究当 $x$ 充分接近 $x_0$ 时函数的变化情况. 例如，考查函数 $y=x\sin\dfrac1x$ 在点 $x=0$ 的附近的变化情况. 我们发现 $\sin\dfrac1x$ 在点 $x=0$ 的邻域内是起伏不定的，但当 $x$ 趋于 $0$ 时，由于
\[
  |y-0|=\left|x\sin\frac1x\right|\le |x|,
\]
$y$ 也随之与 $0$ 无限接近. 因此，我们引入函数极限的定义。

\noindent\textbf{定义 3.1.1}\quad 设函数 $f(x)$ 在 $U_0(x_0,\delta_0)(\delta_0>0)$ 内有定义. 若存在实数 $A$，使得 $\forall\varepsilon>0$，$\exists\delta>0$，当 $x\in U_0(x_0,\delta)$ (即 $0<|x-x_0|<\delta$) 时，有 $|f(x)-A|<\varepsilon$，则称当 $x$ 趋于 $x_0$ 时，函数 $f(x)$ 以 $A$ 为极限 (又称函数 $f(x)$ 在点 $x_0$ 的极限存在，其极限为 $A$)，记为
\[
  \lim_{x\to x_0} f(x)=A
\]
或者 $f(x)\to A\ (x\to x_0)$。

从以上定义中可以看出，$f(x)$ 在 $x_0$ 存在极限 $A$ 与 $f(x)$ 在 $x_0$ 是否有定义无关，即使 $f(x)$ 在 $x_0$ 有定义，$A$ 也不一定等于 $f(x_0)$. 换句话说，在考虑函数极限时，我们只考虑 $x_0$ 附近 $f(x)$ 的变化趋势，而并不关心 $f(x)$ 在 $x_0$ 的取值情况。

再来看 $\lim\limits_{x\to x_0} f(x)=A$ 的几何意义. 任给 $\varepsilon>0$，用平行于 $x$ 轴的直线 $y=A-\varepsilon$ 和 $y=A+\varepsilon$ 作一长条带域. 由定义，存在 $\delta>0$，使得当 $0<|x-x_0|<\delta$ 时，$|f(x)-A|<\varepsilon$. 这就是说，在 $x$ 轴上可以找到一个以 $x_0$ 为中心的开区间 $(x_0-\delta,x_0+\delta)$，使当 $x\in(x_0-\delta,x_0+\delta)$ 且 $x\ne x_0$ 时，点 $P(x,f(x))$ 就必然落在所述的长条阴影带域内 (参见图 3.1.1)。

\noindent\textbf{例 3.1.1}\quad 证明极限
\[
  \lim_{x\to1}\frac{\sqrt{x}-1}{x-1}=\frac12.
\]

\noindent\textbf{分析}\quad 当 $x\ne1$ 时，有
\[
  \left|\frac{\sqrt{x}-1}{x-1}-\frac12\right|
  =\left|\frac{x-1}{2(\sqrt{x}+1)^2}\right|.
\]
由于我们只考虑函数在 $x=1$ 附近的变化情况，因此我们无妨限定 $|x-1|<1$，此时 $(\sqrt{x}+1)^2>1$，从而
\[
  \left|\frac{\sqrt{x}-1}{x-1}-\frac12\right|
  =\frac{|x-1|}{2(\sqrt{x}+1)^2}<\frac12|x-1|.
\]
这样，$\forall\varepsilon>0$，欲使
\[
  \left|\frac{\sqrt{x}-1}{x-1}-\frac12\right|<\varepsilon,
\]
只需取 $\delta=2\varepsilon$。

\noindent\textbf{证明}\quad $\forall\varepsilon>0$，取 $\delta=\min\{1,2\varepsilon\}$，则当 $|x-1|<\delta$ 时，有
\[
  \left|\frac{\sqrt{x}-1}{x-1}-\frac12\right|
  <\frac12|x-1|<\varepsilon.
\]
按函数极限的定义知 $\lim\limits_{x\to1}\dfrac{\sqrt{x}-1}{x-1}=\dfrac12$ 成立。

\noindent\textbf{例 3.1.2}\quad 证明极限 $\lim\limits_{x\to2}x^4=16$。

\noindent\textbf{证明}\quad 由于
\[
  |x^4-16|=|x-2||x^3+2x^2+4x+8|,
\]
限定 $|x-2|<1$ 时，有 $|x|<3$，从而有
\[
  |x^4-16|<65|x-2|.
\]
因此，$\forall\varepsilon>0$，取 $\delta=\min\left\{\dfrac{\varepsilon}{65},1\right\}$，则当 $0<|x-2|<\delta$ 时，有
\[
  |x^4-16|<\varepsilon.
\]
按函数极限的定义知 $\lim\limits_{x\to2}x^4=16$ 成立。

用完全类似的方法可证：对任意的实数 $a$ 和任意的正整数 $n$，有
\[
  \lim_{x\to a}x^n=a^n.
\]
请读者作为练习给出这一结果的证明。

从以上例子的证明可以看出：要证明 $\lim\limits_{x\to x_0}f(x)=A$，我们应该先限定 $x$ 在 $x_0$ 的某个邻域内，这样容易从 $|f(x)-A|$ 中得到所要的关于 $|x-x_0|$ 的不等式. 而预先限定 $x$ 的变化范围，可以通过给定某个 $\delta_0>0$ 得到。

\noindent\textbf{例 3.1.3}\quad 证明极限 $\lim\limits_{x\to0}x\sin\dfrac1x=0$。

\noindent\textbf{证明}\quad $\forall\varepsilon>0$，取 $\delta=\varepsilon$，当 $x\in U_0(0,\delta)$ 时，有
\[
  \left|x\sin\frac1x-0\right|\le |x|<\varepsilon,
\]
按函数极限的定义知 $\lim\limits_{x\to0}x\sin\dfrac1x=0$ 成立。

\subsection{函数极限的性质}

与序列极限的性质一样，函数极限也具有相应的性质. 但是由于函数极限的存在与否只与 $x_0$ 附近的函数值有关，因此一些相应的结论也只能在 $x_0$ 的某个邻域内成立. 例如，对于``改变序列的有限多项不改变序列的敛散性''这一性质而言，我们有以下相应的结论：设 $f(x)$ 在 $x_0$ 的某个邻域 $U_0(x_0,\delta_0)$ ($\delta_0>0$) 有定义，则对任何 $\delta_0>\delta_1>0$，改变 $f(x)$ 在 $U_0(x_0,\delta_1)$ 外面的函数值，不影响 $f(x)$ 在 $x_0$ 处的敛散性。

\noindent\textbf{定理 3.1.1 (唯一性和有界性)}\quad 设极限 $\lim\limits_{x\to x_0}f(x)$ 存在，则
\begin{enumerate}
  \item 该极限值是唯一的；
  \item $\exists\delta_0>0$，使得 $f(x)$ 在 $U_0(x_0,\delta_0)$ 有界。
\end{enumerate}

\noindent\textbf{证明}\quad (1) 设 $\lim\limits_{x\to x_0}f(x)=A$，$\lim\limits_{x\to x_0}f(x)=B$，则由函数极限的定义，$\forall\varepsilon>0$，$\exists\delta>0$，当 $0<|x-x_0|<\delta$ 时，有
\[
  |f(x)-A|<\frac{\varepsilon}{2},\qquad |f(x)-B|<\frac{\varepsilon}{2},
\]
从而有
\[
  |A-B|<|f(x)-A|+|f(x)-B|<\varepsilon.
\]
由于 $A,B$ 为两个常数，它们的差又能小于任给的 $\varepsilon>0$，因此 $A$ 必须等于 $B$。

(2) 设 $\lim\limits_{x\to x_0}f(x)=A$. 对 $\varepsilon_0=1$，$\exists\delta_0>0$，当 $0<|x-x_0|<\delta_0$ 时，有 $|f(x)-A|<1$，从而有
\[
  |f(x)|\le |A|+|f(x)-A|<1+|A|,\qquad \forall x\in U_0(x_0,\delta_0),
\]
按有界函数的定义知 $f(x)$ 在 $U_0(x_0,\delta_0)$ 有界。

\noindent\textbf{定理 3.1.2 (保序性)}\quad 设 $\lim\limits_{x\to x_0}f(x)=A$，$\lim\limits_{x\to x_0}g(x)=B$，则
\begin{enumerate}
  \item 若 $\exists\delta_0>0$，使得当 $x\in U_0(x_0,\delta_0)$ 时，有 $f(x)\le g(x)$，则 $A\le B$；
  \item 若 $A<B$，则对任何 $C\in(A,B)$，$\exists\delta_C>0$，使得当 $x\in U_0(x_0,\delta_C)$ 时，有 $f(x)<C<g(x)$。
\end{enumerate}
定理的证明留给读者作为练习。

\noindent\textbf{定理 3.1.3 (四则运算)}\quad 设 $\lim\limits_{x\to x_0}f(x)=A$，$\lim\limits_{x\to x_0}g(x)=B$，则有
\begin{enumerate}
  \item $\lim\limits_{x\to x_0}(f(x)\pm g(x))=A\pm B$；
  \item $\lim\limits_{x\to x_0}f(x)\cdot g(x)=A\cdot B$；
  \item $\lim\limits_{x\to x_0}\dfrac{f(x)}{g(x)}=\dfrac{A}{B}$，这里假定 $B\ne0$。
\end{enumerate}

\noindent\textbf{证明}\quad 我们只证 (3)，其余的请读者自己证明。

在 $x_0$ 的附近，我们有
\[
  \left|\frac{f(x)}{g(x)}-\frac{A}{B}\right|
  =\frac{|Bf(x)-Ag(x)|}{|Bg(x)|}
  =\frac{|B(f(x)-A)-A(g(x)-B)|}{|B||g(x)|}
\]
\[
  \le \frac{1}{|B||g(x)|}\bigl(|B||f(x)-A|+|A||g(x)-B|\bigr).
\]
由 $\lim\limits_{x\to x_0}g(x)=B\ne0$ 知，$\exists\delta_0>0$，使得当 $x\in U_0(x_0,\delta_0)$ 时，有
\[
  |g(x)-B|<\frac{|B|}{2},
\]
从而
\[
  |g(x)|\ge |B|-|g(x)-B|\ge |B|-\frac{|B|}{2}=\frac{|B|}{2}.
\]
这样，与前面的不等式相结合，我们有
\[
  \left|\frac{f(x)}{g(x)}-\frac{A}{B}\right|
  \le \frac{2}{|B|}|f(x)-A|+2\frac{|A|}{|B|^2}|g(x)-B|
\]
对一切的 $x\in U_0(x_0,\delta_0)$ 成立。

$\forall\varepsilon>0$，由 $\lim\limits_{x\to x_0}g(x)=B$ 知，$\exists\delta_1>0$，使得当 $x\in U_0(x_0,\delta_1)$ 时，有
\[
  |g(x)-B|<\frac{B^2}{4(|A|+1)}\varepsilon;
\]
再由 $\lim\limits_{x\to x_0}f(x)=A$ 知，$\exists\delta_2>0$，使得当 $x\in U_0(x_0,\delta_2)$ 时，有
\[
  |f(x)-A|<\frac{|B|}{4}\varepsilon.
\]

现在取 $\delta=\min\{\delta_0,\delta_1,\delta_2\}$，则当 $x\in U_0(x_0,\delta)$ 时，有
\[
  \left|\frac{f(x)}{g(x)}-\frac{A}{B}\right|
  <\frac{\varepsilon}{2}+\frac{\varepsilon}{2}=\varepsilon.
\]
按函数极限的定义知 (3) 成立。

\noindent\textbf{例 3.1.4}\quad 求极限
\[
  \lim_{x\to2}\frac{2x^3+3x-1}{x^3+3}.
\]

\noindent\textbf{解}\quad 因 $x\to2$ 时，$2x^3+3x-1$ 和 $x^3+3$ 都有极限，而且
\[
  \lim_{x\to2}(2x^3+3x-1)=2\times2^3+3\times2-1=21,
\]
\[
  \lim_{x\to2}(x^3+3)=2^3+3=11\ne0,
\]
所以
\[
  \lim_{x\to2}\frac{2x^3+3x-1}{x^3+3}
  =\frac{\lim\limits_{x\to2}(2x^3+3x-1)}
  {\lim\limits_{x\to2}(x^3+3)}
  =\frac{21}{11}.
\]

\noindent\textbf{例 3.1.5}\quad 求
\[
  \lim_{x\to1}\frac{x+x^2+\cdots+x^n-n}{x-1},
\]
其中 $n$ 是一个给定的正整数。

\noindent\textbf{解}\quad 将所求极限式恒等变形，有
\[
\begin{aligned}
  \lim_{x\to1}\frac{x+x^2+\cdots+x^n-n}{x-1}
  &=\lim_{x\to1}\left(\frac{x-1}{x-1}+\frac{x^2-1}{x-1}+\cdots+\frac{x^n-1}{x-1}\right)\\
  &=\lim_{x\to1}\left[1+(1+x)+\cdots+(x^{n-1}+x^{n-2}+\cdots+1)\right]\\
  &=1+2+\cdots+n=\frac12 n(n+1).
\end{aligned}
\]

下面我们来讨论复合函数求极限的问题，有下述结论：

\noindent\textbf{定理 3.1.4}\quad 设函数 $f(u)$ 在 $U_0(u_0,\delta_1)(\delta_1>0)$ 有定义且 $\lim\limits_{u\to u_0}f(u)=A$；$u=g(x)$ 在 $U_0(x_0,\delta_0)(\delta_0>0)$ 有定义，当 $x\in U_0(x_0,\delta_0)$ 时有 $g(x)\in U_0(u_0,\delta_1)$ 且 $\lim\limits_{x\to x_0}g(x)=u_0$，则有
\[
  \lim_{x\to x_0}f(g(x))=A.
\]

\noindent\textbf{证明}\quad 由 $\lim\limits_{u\to u_0}f(u)=A$ 知，$\forall\varepsilon>0$，$\exists 0<\eta<\delta_1$，使得当 $u\in U_0(u_0,\eta)$ 时，有
\[
  |f(u)-A|<\varepsilon.
\]
对于 $\eta>0$，由 $\lim\limits_{x\to x_0}g(x)=u_0$，$\exists 0<\delta<\delta_0$，使得当 $x\in U_0(x_0,\delta)$ 时，有
\[
  |g(x)-u_0|<\eta.
\]
注意到 $\forall x\in U_0(x_0,\delta_0)$ 有 $g(x)\ne u_0$，因此当 $x\in U_0(x_0,\delta)$ 时，有 $g(x)\in U_0(u_0,\eta)$，因此有
\[
  |f(g(x))-A|<\varepsilon.
\]
这就证明了 $\lim\limits_{x\to x_0}f(g(x))=A$. 定理证毕。

\noindent\textbf{例 3.1.6}\quad 求极限
\[
  \lim_{x\to0}\frac{\sqrt{x^2+x+1}-1}{x^2+x}.
\]

\noindent\textbf{解}\quad 取 $f(u)=\dfrac{\sqrt{u}-1}{u-1}$，其中 $u=x^2+x+1$，则 $f(u)$ 在 $u_0=1$ 的某个邻域，$g(x)$ 在 $x_0=0$ 的某个邻域内满足定理 3.1.4 的所有条件. 由例 3.1.1 知 $\lim\limits_{u\to1}f(u)=\dfrac12$. 显然有 $\lim\limits_{x\to0}g(x)=\lim\limits_{x\to0}(x^2+x+1)=1$. 由定理 3.1.4 知
\[
  \lim_{x\to0}\frac{\sqrt{x^2+x+1}-1}{x^2+x}
  =\lim_{x\to0}f(g(x))=\frac12.
\]

对于定理 3.1.4 中的条件，我们要求当 $x\in U_0(x_0,\delta_0)$ 时，有 $g(x)\in U_0(u_0,\delta_1)$. 该条件在复合函数求极限中是必要的. 换句话说，若假定 $g(x)\in U(u_0,\delta_1)$，则定理 3.1.4 的结论可能不真. 事实上，我们取
\[
  f(u)=
  \begin{cases}
    0, & u=0,\\
    1, & u\ne0,
  \end{cases}
  \qquad
  g(x)=
  \begin{cases}
    1, & x=\dfrac1n,\\
    0, & \text{其他},
  \end{cases}
\]
则
\[
  f(g(x))=
  \begin{cases}
    1, & x=\dfrac1n,\\
    0, & \text{其他}.
  \end{cases}
\]
现取 $x_0=0$，则 $\lim\limits_{x\to0}g(x)=0$. 注意到 $\lim\limits_{u\to0}f(u)=1$，但是 $\lim\limits_{x\to0}f(g(x))$ 不存在. 这是因为在 $x_0=0$ 的任意小的邻域内 $f(g(x))$ 既有取值为 $1$ 的点，又有取值为 $0$ 的点，故其极限不存在。

\noindent\textbf{定理 3.1.5 (夹逼收敛原理)}\quad 若 $\lim\limits_{x\to x_0}f(x)=\lim\limits_{x\to x_0}h(x)=A$，而且存在 $\delta_0>0$，使得 $f(x)\le g(x)\le h(x)$ 对一切 $x\in U_0(x_0,\delta)$ 成立，则有
\[
  \lim_{x\to x_0}g(x)=A.
\]

此定理的证明与序列的逼收敛原理类似，故从略。

最后，我们要重复指出的是，函数极限存在时，它所具有的性质只是在局部成立. 例如，设 $f(x)=\dfrac1x$，$x\in(0,+\infty)$. 对任意的 $x_0\in(0,+\infty)$，有
\[
  \lim_{x\to x_0}f(x)=\lim_{x\to x_0}\frac1x=\frac1{x_0}.
\]
由此可以导出 $f(x)$ 在 $x_0$ 附近有界，事实上，我们有
\[
  |f(x)|<\frac{2}{x_0},\qquad \forall x\in\left(x_0-\frac{x_0}{2},x_0+\frac{x_0}{2}\right),
\]
但 $f(x)$ 显然在 $(0,+\infty)$ 上无界。

\subsection{函数极限概念的推广}

在实际应用中，前面所引入的函数极限的概念是不够的，例如对闭区间上有定义的函数，就无法讨论其在端点处的极限问题. 因此我们就必须拓广极限的概念，以便处理其他形式的极限问题。

\noindent\textbf{1. 单侧极限}

所谓\textbf{单侧极限}是考虑函数在某一点的一侧的变化情况. 记
\[
  U^+(x_0,\delta)=\{x\in\mathbb{R}:x_0\le x<x_0+\delta\}\qquad (x_0\text{ 的右邻域});
\]
\[
  U^-(x_0,\delta)=\{x\in\mathbb{R}:x_0-\delta<x\le x_0\}\qquad (x_0\text{ 的左邻域});
\]
\[
  U_0^+(x_0,\delta)=U^+(x_0,\delta)\setminus\{x_0\}\qquad (x_0\text{ 的右空心邻域});
\]
\[
  U_0^-(x_0,\delta)=U^-(x_0,\delta)\setminus\{x_0\}\qquad (x_0\text{ 的左空心邻域}).
\]

\noindent\textbf{定义 3.1.2}\quad 设 $f(x)$ 在 $U_0^+(x_0,\delta_0)(\delta_0>0)$ 上有定义. 如果存在实数 $A$，对 $\forall\varepsilon>0$，$\exists\delta>0$，使得当 $x\in U_0^+(x_0,\delta)$ (即 $0<x-x_0<\delta$) 时，有 $|f(x)-A|<\varepsilon$，则称 $f(x)$ 在点 $x_0$ 的\textbf{右极限存在}，而称 $A$ 为 $f(x)$ 在点 $x_0$ 的\textbf{右极限}，记为
\[
  \lim_{x\to x_0+0}f(x)=A
\]
或者 $f(x_0+0)=A$。

类似地可定义 $f(x)$ 在点 $x_0$ 的\textbf{左极限} $\lim\limits_{x\to x_0-0}f(x)$ 或者 $f(x_0-0)$. 有关右极限的几何解释参见图 3.1.2。

\noindent\textbf{定理 3.1.6}\quad 函数 $f(x)$ 在点 $x_0$ 极限存在的充分必要条件是 $f(x)$ 在点 $x_0$ 的左、右极限都存在且相等。

\noindent\textbf{证明}\quad 必要性\quad 设 $\lim\limits_{x\to x_0}f(x)=A$ 存在，则 $\forall\varepsilon>0$，$\exists\delta>0$，使得当 $0<|x-x_0|<\delta$ 时，有
\[
  |f(x)-A|<\varepsilon.
\]
特别地，当 $0<x-x_0<\delta$ 或 $0<x_0-x<\delta$ 时，亦有 $|f(x)-A|<\varepsilon$，故 $\lim\limits_{x\to x_0+0}f(x)$ 和 $\lim\limits_{x\to x_0-0}f(x)$ 都存在，而且
\[
  \lim_{x\to x_0+0}f(x)=\lim_{x\to x_0-0}f(x)=A.
\]

充分性\quad 设 $\lim\limits_{x\to x_0+0}f(x)$ 和 $\lim\limits_{x\to x_0-0}f(x)$ 都存在，而且
\[
  \lim_{x\to x_0+0}f(x)=\lim_{x\to x_0-0}f(x)=A,
\]
则 $\forall\varepsilon>0$，由 $\lim\limits_{x\to x_0+0}f(x)=A$ 知，$\exists\delta_1>0$，使得当 $0<x-x_0<\delta_1$ 时，有 $|f(x)-A|<\varepsilon$；又由 $\lim\limits_{x\to x_0-0}f(x)=A$ 知，$\exists\delta_2>0$，使得当 $0<x_0-x<\delta_2$ 时，有 $|f(x)-A|<\varepsilon$。

令 $\delta=\min\{\delta_1,\delta_2\}$，则当 $0<|x-x_0|<\delta$ 时，有 $|f(x)-A|<\varepsilon$. 故 $\lim\limits_{x\to x_0}f(x)$ 存在，而且 $\lim\limits_{x\to x_0}f(x)=A$。

\noindent\textbf{例 3.1.7}\quad 证明取整函数 $f(x)=[x]$ 在 $x=3$ 的极限不存在。

\noindent\textbf{证明}\quad 由取整函数的定义可证
\[
  f(3+0)=3>2=f(3-0),
\]
因此，由定理 3.1.5 知，取整函数 $f(x)=[x]$ 在 $x=3$ 的极限不存在。

\noindent\textbf{2. 自变量趋向无穷大时的极限}

如果将序列看成是定义在 $\mathbb{N}$ 上的函数的话，则序列极限就是自变量 $n\to+\infty$ 时函数的极限，类似地也可以考虑在实数集 $\mathbb{R}$ 上定义的函数当 $x$ 趋于 $\infty$ 时的极限问题。

称集合 $\{x:|x|>h\}$ ($h>0$) 为 $\infty$ 的邻域，记为 $U(\infty,h)$ 或 $U(\infty)$；称 $\{x:h<x<+\infty\}$ 与 $\{x:-\infty<x<-h\}$ 为 $\infty$ 的单侧邻域，分别记为 $U^+(\infty,h)(U^+(\infty)\text{ 或 }U(+\infty))$ 和 $U^-(\infty,h)(U^-(\infty)\text{ 或 }U(-\infty))$。

\noindent\textbf{定义 3.1.3}\quad 设函数 $f(x)$ 在 $U^+(\infty)$ 上有定义. 若存在实数 $A$，使得 $\forall\varepsilon>0$，$\exists X\in U^+(\infty)$，当 $x>X$ 时，有 $|f(x)-A|<\varepsilon$，则称当 $x$ 趋于 $+\infty$ 时 $f(x)$ 的极限存在，其极限为 $A$，记为 $\lim\limits_{x\to+\infty}f(x)=A$ 或者 $f(x)\to A\ (x\to+\infty)$。

类似地可以定义 $\lim\limits_{x\to-\infty}f(x)=A$ 和 $\lim\limits_{x\to\infty}f(x)=A$. 有关 $x\to+\infty$ 时 $f(x)$ 的极限的几何解释参见图 3.1.3。

对这种类型的极限，也有类似定理 3.1.6 结论，另外，这类极限也有四则运算及复合函数相应的结果. 请读者作为练习给出它们的表述和证明。

\noindent\textbf{例 3.1.8}\quad 证明极限
\[
  \lim_{x\to\infty}\frac{3x^2+2x-1}{x^2-2}=3.
\]

\noindent\textbf{证明}\quad 由于
\[
  \left|\frac{3x^2+2x-1}{x^2-2}-3\right|
  =\left|\frac{2x+5}{x^2-2}\right|
  \le\frac{3|x|}{3x^2/4}=\frac{4}{|x|}
\]
对一切的 $|x|>5$ 成立，所以 $\forall\varepsilon>0$，取 $X=\max\{5,4/\varepsilon\}$，则当 $|x|>X$ 时，有
\[
  \left|\frac{3x^2+2x-1}{x^2-2}-3\right|\le\frac4{|x|}<\varepsilon.
\]
按极限定义知 $\lim\limits_{x\to\infty}\dfrac{3x^2+2x-1}{x^2-2}=3$ 成立。

\noindent\textbf{例 3.1.9}\quad 设 $f(x)=\sqrt{4x^2+2x+3}$，试求下列极限：
\[
  (1)\ \lim_{x\to+\infty}\frac{f(x)}{x};\qquad
  (2)\ \lim_{x\to-\infty}\frac{f(x)}{x};
\]
\[
  (3)\ \lim_{x\to+\infty}(f(x)-2x);\qquad
  (4)\ \lim_{x\to-\infty}(f(x)+2x).
\]

\noindent\textbf{解}\quad (1)
\[
  \lim_{x\to+\infty}\frac{f(x)}{x}
  =\lim_{x\to+\infty}\sqrt{4+\frac2x+\frac3{x^2}}=2;
\]
(2)
\[
  \lim_{x\to-\infty}\frac{f(x)}{x}
  =\lim_{t\to+\infty}\frac{\sqrt{4t^2-2t+3}}{-t}=-2;
\]
(3)
\[
  \lim_{x\to+\infty}(f(x)-2x)
  =\lim_{x\to+\infty}\frac{4x^2+2x+3-4x^2}{\sqrt{4x^2+2x+3}+2x}
  =\frac12;
\]
(4)
\[
\begin{aligned}
  \lim_{x\to-\infty}(f(x)+2x)
  &=\lim_{x\to-\infty}\frac{4x^2+2x+3-4x^2}{\sqrt{4x^2+2x+3}-2x}\\
  &=\lim_{t\to+\infty}\frac{-2t+3}{\sqrt{4t^2-2t+3}+2t}
  =-\frac12.
\end{aligned}
\]

\noindent\textbf{3. 广义极限}

当函数 $f(x)$ 在点 $x_0$ 的极限不存在时，这时 $f(x)$ 在 $x_0$ 附近的变化情况非常复杂. 但所有在点 $x_0$ 不存在极限的函数中，有一类变化还是具有规律的. 如函数 $f(x)=\dfrac1x$，当 $x\to0$ 时，函数的绝对值将目标一致地无限变大。

\noindent\textbf{定义 3.1.4}\quad 设 $f(x)$ 在 $U_0(x_0,\delta_0)(\delta_0>0)$ 上有定义. 若 $\forall M>0$，$\exists\delta>0$，使得当 $0<|x-x_0|<\delta$ 时，有 $f(x)>M$，则称当 $x$ 趋于 $x_0$ 时，$f(x)$ 趋于 $+\infty$，或称当 $x$ 趋于 $x_0$ 时，$f(x)$ 的广义极限为 $+\infty$，并记为 $\lim\limits_{x\to x_0}f(x)=+\infty$ 或 $f(x)\to+\infty\ (x\to x_0)$. 此时，也称 $f(x)$ 为当 $x$ 趋于 $x_0$ 时的\textbf{正无穷大量}。

类似地可以定义 $\lim\limits_{x\to x_0}f(x)=-\infty$ (称 $f(x)$ 为当 $x$ 趋于 $x_0$ 时的\textbf{负无穷大量})，$\lim\limits_{x\to x_0}f(x)=\infty$ (称 $f(x)$ 为当 $x$ 趋于 $x_0$ 时的\textbf{无穷大量})，$\lim\limits_{x\to x_0+0}f(x)=+\infty$ (称 $f(x)$ 为当 $x$ 从右侧趋于 $x_0$ 时的\textbf{正无穷大量})，$\lim\limits_{x\to+\infty}f(x)=+\infty$ (称 $f(x)$ 为当 $x$ 趋于 $+\infty$ 时的\textbf{正无穷大量}) 等。

注意无穷大量并不是函数值可以取到任意大的值的量. 例如，函数 $f(x)=x^2\sin x$ 当 $x\to\infty$ 时可以取到任意大的值，但是它并不是当 $x\to\infty$ 时的无穷大量。

对于函数极限而言，自变量 $x$ 共有六种可能的变化情况，即
\[
  x\to x_0;\quad x\to x_0+0;\quad x\to x_0-0;\quad x\to\infty;\quad x\to\pm\infty.
\]
对于自变量 $x$ 的每一种趋向，函数又有四种不同的趋向，即
\[
  f(x)\to A;\quad f(x)\to\pm\infty;\quad f(x)\to\infty.
\]
因此在研究函数极限时，就有 24 种可能的情形要研究. 对于
\[
  \lim_{x\to x_0\pm0}f(x)=A,\qquad
  \lim_{x\to\infty}f(x)=A,\qquad
  \lim_{x\to\pm\infty}f(x)=A
\]
这五种极限，也相应地有上一小节中所介绍的所有性质；而对于其他类型的极限，在不发生 $(\pm\infty)+(\mp\infty)$，$0\cdot\infty$，$\dfrac00$ 和 $\dfrac{\infty}{\infty}$ 这些类型 (称为不定式) 情况下，也有相应的四则运算公式。

\noindent\textbf{例 3.1.10}\quad 求极限
\[
  \lim_{x\to\infty}\frac{5-2x^2}{3x^2+x-2}.
\]

\noindent\textbf{解}\quad 利用极限的四则运算，可得
\[
\begin{aligned}
  \lim_{x\to\infty}\frac{5-2x^2}{3x^2+x-2}
  &=\lim_{x\to\infty}\frac{\dfrac5{x^2}-2}{3+\dfrac1x-\dfrac2{x^2}}\\
  &=\frac{\lim\limits_{x\to\infty}\left(\dfrac5{x^2}-2\right)}
  {\lim\limits_{x\to\infty}\left(3+\dfrac1x-\dfrac2{x^2}\right)}
  =-\frac23.
\end{aligned}
\]

\noindent\textbf{例 3.1.11}\quad 设 $f(x)=\dfrac1{x-1}$，求 $f(1-0)$，$f(1+0)$ 和 $\lim\limits_{x\to1}f(x)$。

\noindent\textbf{解}\quad $f(1-0)=-\infty$，$f(1+0)=+\infty$ 和 $\lim\limits_{x\to1}f(x)=\infty$。

\subsection{序列极限与函数极限的关系}

以 $x$ 趋于 $x_0$ 的情形为例，我们先来给出 $f(x)$ 不趋于 $A$ 的精确描述. 如同序列情形，$f(x)$ 不趋于 $A$，则必定 $\exists\varepsilon_0>0$，在 $x_0$ 的附近有充分接近 $x_0$ 的点 $x'$，使得 $|f(x')-A|\ge\varepsilon_0$. 读者应该充分理解``有充分接近 $x_0$ 的点 $x'$''的精确意义，此即是说，在 $x_0$ 的任何邻域内都存在异于 $x_0$ 的点 $x'$，使得 $|f(x')-A|\ge\varepsilon_0$. 因此，用肯定的语气来叙述 $\lim\limits_{x\to x_0}f(x)\ne A$，即为
\[
  \exists\varepsilon_0>0,\ \forall\delta>0,\ \exists x'\in U_0(x_0,\delta),\ \text{使得 } |f(x')-A|\ge\varepsilon_0.
\]
我们这里须特别指出的是，初学者容易犯的一个错误是，将其叙述成：
\[
  \exists\varepsilon_0>0,\ \forall\delta>0,\ \text{当 }x\in U_0(x_0,\delta)\text{ 时，有 } |f(x)-A|\ge\varepsilon_0.
\]
当然，满足以上条件一定有 $x$ 趋于 $x_0$ 时 $f(x)$ 不趋于 $A$，但并不是所有的 $x$ 趋于 $x_0$ 时 $f(x)$ 不以 $A$ 为极限的都有以上情形发生。

下面我们来讨论序列极限与函数极限的关系问题. 事实上，我们有

\noindent\textbf{定理 3.1.7}\quad 设 $f(x)$ 在 $U_0(x_0,\delta_0)(\delta_0>0)$ 上有定义，则 $\lim\limits_{x\to x_0}f(x)=A$ 成立的充分必要条件是：对于 $U_0(x_0,\delta_0)$ 内任意收敛于 $x_0$ 的序列 $\{x_n\}$，都有
\[
  \lim_{n\to\infty}f(x_n)=A.
\]

\noindent\textbf{证明}\quad 必要性\quad 在 $U_0(x_0,\delta_0)$ 中任取一收敛于 $x_0$ 的序列 $\{x_n\}$，即 $\lim\limits_{n\to\infty}x_n=x_0$，要证 $\lim\limits_{n\to\infty}f(x_n)=A$。

$\forall\varepsilon>0$，由 $\lim\limits_{x\to x_0}f(x)=A$ 知，$\exists\delta>0$，使得当 $0<|x-x_0|<\delta$ 时，有
\[
  |f(x)-A|<\varepsilon.
\]
对于上述 $\delta>0$，由 $x_n\to x_0\ (n\to\infty)$ 知，$\exists N$，使得当 $n>N$ 时，有
\[
  0<|x_n-x_0|<\delta.
\]
于是，当 $n>N$ 时，有 $|f(x_n)-A|<\varepsilon$，即 $\lim\limits_{n\to\infty}f(x_n)=A$。

充分性\quad 如果不然，即当 $x\to x_0$ 时 $f(x)$ 不以 $A$ 为极限，则 $\exists\varepsilon_0>0$，对 $\forall\delta>0$，$\exists x_\delta\in U_0(x_0,\delta)$，使得
\[
  |f(x_\delta)-A|\ge\varepsilon_0.
\]
这样，分别取 $\delta=\dfrac{\delta_0}{n}\ (n=1,2,\cdots)$，$\exists x_n\in U_0\left(x_0,\dfrac{\delta_0}{n}\right)$，使得
\[
  |f(x_n)-A|\ge\varepsilon_0.
\]
对于这样得到的序列 $\{x_n\}$，显然有 $x_n\to x_0\ (n\to\infty)$，且 $x_n\in U_0(x_0,\delta_0)$，从而 $\lim\limits_{n\to\infty}f(x_n)=A$，但这与 $|f(x_n)-A|\ge\varepsilon_0$ 对一切正整数 $n$ 成立相矛盾。

\noindent\textbf{例 3.1.12}\quad 试证明当 $x\to0$ 时函数 $f(x)=\sin\dfrac1x\ (x\ne0)$ 的极限不存在。

\noindent\textbf{证明}\quad 取
\[
  x_n=\frac1{n\pi},\qquad n=1,2,\cdots,
\]
则有 $x_n\to0\ (n\to\infty)$，且
\[
  f(x_n)=\sin n\pi\equiv0\to0\qquad (n\to\infty);
\]
再取
\[
  \widetilde{x}_n=\frac1{2n\pi+\dfrac{\pi}{2}},\qquad n=1,2,\cdots,
\]
则有 $\widetilde{x}_n\to0\ (n\to\infty)$，且
\[
  f(\widetilde{x}_n)=\sin\left(2n\pi+\frac{\pi}{2}\right)\equiv1\to1\qquad (n\to\infty).
\]
这样，由定理 3.1.7 知 $\lim\limits_{x\to0}\sin\dfrac1x$ 不存在。

值得注意的是，对于此例，虽然当 $x\to0$ 时 $\sin\dfrac1x$ 不趋向于 $0$，但找不到 $\varepsilon_0>0$，使得对任何 $\delta>0$，都有 $\left|\sin\dfrac1x-0\right|\ge\varepsilon_0$ 在 $U_0(x_0,\delta)$ 上成立。

\subsection{极限存在性定理和两个重要极限}

\noindent\textbf{1. 极限存在性定理}

如同序列极限一样，如果函数单调，则它必存在广义极限。

\noindent\textbf{定理 3.1.8}\quad 设函数 $f(x)$ 在 $U_0^+(x_0,\delta_0)$ 内有定义，若 $f(x)$ 在 $U_0^+(x_0,\delta_0)$ 上内单调上升，则
\[
  \lim_{x\to x_0+0}f(x)=\inf\{f(x):x\in U_0^+(x_0,\delta_0)\};
\]
若 $f(x)$ 在 $U_0^+(x_0,\delta_0)$ 内单调下降，则
\[
  \lim_{x\to x_0+0}f(x)=\sup\{f(x):x\in U_0^+(x_0,\delta_0)\}.
\]

\noindent\textbf{证明}\quad 我们只给出定理第一部分的证明. 记
\[
  A=\inf\{f(x):x\in U_0^+(x_0,\delta_0)\},
\]
则 $A$ 为有限数或 $-\infty$。

先考虑 $A$ 为有限数的情形. 对 $\forall\varepsilon>0$，由下确界的性质知，$\exists x_1\in U_0^+(x_0,\delta_0)$，使得 $A\le f(x_1)<A+\varepsilon$. 记 $\delta=x_1-x_0$，则 $\delta>0$，而且当 $x\in U_0^+(x_0,\delta)$ 时，由于 $f(x)$ 在 $U_0^+(x_0,\delta_0)$ 内单调上升，故有 $A\le f(x)\le f(x_1)<A+\varepsilon$，即 $|f(x)-A|<\varepsilon$. 这就证明了，此时 $\lim\limits_{x\to x_0+0}f(x)=A$ 成立。

再考虑 $A=-\infty$ 的情形. 对 $\forall M>0$，由下确界的定义知，$\exists x_1\in U_0^+(x_0,\delta_0)$ 使得 $f(x_1)<-M$. 记 $\delta=x_1-x_0$，则 $\delta>0$，而且当 $x\in U_0^+(x_0,\delta)$ 时，由于 $f(x)$ 在 $U_0^+(x_0,\delta_0)$ 内单调上升，故有 $f(x)\le f(x_1)<-M$. 这就证明了，此时亦有 $\lim\limits_{x\to x_0+0}f(x)=A$ 成立。

对于函数极限，有下面的柯西收敛准则：

\noindent\textbf{定理 3.1.9}\quad 设 $f(x)$ 在 $U_0(x_0,\delta_0)$ 内有定义，则 $\lim\limits_{x\to x_0}f(x)$ 存在的充分必要条件是：$\forall\varepsilon>0$，$\exists\delta>0$，当 $x',x''\in U_0(x_0,\delta)$ 时，有
\[
  |f(x')-f(x'')|<\varepsilon.
\]

\noindent\textbf{证明}\quad 必要性\quad 设 $\lim\limits_{x\to x_0}f(x)=A$ 存在，则 $\forall\varepsilon>0$，$\exists\delta>0$，当 $x\in U_0(x_0,\delta)$ 时，有 $|f(x)-A|<\varepsilon/2$. 于是，当 $x',x''\in U_0(x_0,\delta)$ 时，有
\[
  |f(x')-f(x'')|\le |f(x')-A|+|A-f(x'')|<\frac{\varepsilon}{2}+\frac{\varepsilon}{2}=\varepsilon.
\]

充分性\quad 在 $U_0(x_0,\delta_0)$ 内任意选取一个收敛于 $x_0$ 的序列 $\{x_n\}$. $\forall\varepsilon>0$，由充分性假定知，$\exists\delta>0$，使得
\[
  |f(x')-f(x'')|<\varepsilon/2,\qquad \forall x',x''\in U_0(x_0,\delta).
\]
而对这一 $\delta>0$，由 $\lim\limits_{n\to\infty}x_n=x_0$ 知，存在 $N$，使得当 $n>N$ 时，有 $|x_n-x_0|<\delta$. 于是，对 $\forall m,n>N$，有 $x_n,x_m\in U_0(x_0,\delta)$，从而有
\[
  |f(x_n)-f(x_m)|<\varepsilon/2.
\]
这表明序列 $\{f(x_n)\}$ 是一个柯西序列，因此由序列的柯西收敛准则知，存在实数 $A$，使得 $\lim\limits_{n\to\infty}f(x_n)=A$。

下证 $\lim\limits_{x\to x_0}f(x)=A$. 由前面所证知，$\forall\varepsilon>0$，存在 $\delta>0$ 和正整数 $N$，使得
\[
  |f(x)-f(x_n)|<\varepsilon/2,\qquad \forall x\in U_0(x_0,\delta),\ \forall n>N.
\]
于是，有
\[
  |f(x)-A|=\lim_{n\to\infty}|f(x)-f(x_n)|\le\varepsilon/2<\varepsilon,\qquad \forall x\in U_0(x_0,\delta).
\]
这就证明了 $\lim\limits_{x\to x_0}f(x)=A$ 成立。

\noindent\textbf{注}\quad 对于其他类型极限也有相应的单调收敛性定理和柯西收敛准则，请读者作为练习给出这些结果的详细表述。

\noindent\textbf{例 3.1.13}\quad 证明极限
\[
  \lim_{x\to+\infty}\frac{x-1}{x+1}\cos\frac{2\pi x}{3}
\]
不存在。

\noindent\textbf{证明}\quad 令 $\varepsilon_0=\dfrac12$，对任意的 $X>0$ (不妨设 $X>3$)，取正整数 $n_0$，使得
\[
  x'=3n_0>X,\qquad x''=3n_0+\frac34>X,
\]
则有
\[
  |f(x')-f(x'')|=\left|\frac{3n_0-1}{3n_0+1}\right|>\frac12=\varepsilon_0.
\]
故由柯西收敛准则知，$\lim\limits_{x\to+\infty}\dfrac{x-1}{x+1}\cos\dfrac{2\pi x}{3}$ 不存在。

\noindent\textbf{2. 两个重要极限}

\noindent\textbf{第一个重要极限}\quad
\[
  \lim_{x\to0}\frac{\sin x}{x}=1.
\]

\noindent\textbf{分析}\quad 如图 3.1.4 所示，在单位圆盘 $D=\{(\xi,\eta):\xi^2+\eta^2\le1\}$ 上，$x$ 是圆心角 $\angle AOB$，以弧度计，即它恰好等于 $\overset{\frown}{AB}$ 的弧长，而 $\sin x=\overline{BC}$ 是弦长 $\overline{BB'}$ 之半. 因此
\[
  \frac{\overline{BB'}}{\overset{\frown}{BB'}}
  =\frac{2\sin x}{2x}=\frac{\sin x}{x}\to1\qquad (x\to0),
\]
即圆心角趋于 $0$ 时，对应的弦长与弧长之比趋于 $1$。

\noindent\textbf{证明}\quad 为了给出这一重要极限的证明，我们先证明一个重要的不等式：
\[
  \sin x<x<\tan x,\qquad \forall x\in\left(0,\frac{\pi}{2}\right).
\]
事实上，如图 3.1.4 所示，我们有
\[
  \triangle AOB\text{ 的面积}<\text{扇形 }AOB\text{ 的面积}<\triangle AOD\text{ 的面积},
\]
即
\[
  \frac12\sin x<\frac12x<\frac12\tan x.
\]
这样，我们就证明了所述的不等式。

利用这一不等式，容易导出，对一切的 $0\ne x\in\mathbb{R}$，有
\[
  |\sin x|<|x|.
\]
由此可以导出，对 $\forall a\in\mathbb{R}$，有
\[
  |\cos x-\cos a|
  =\left|2\sin\frac{x+a}{2}\sin\frac{x-a}{2}\right|
  \le 2\left|\sin\frac{x-a}{2}\right|\le |x-a|,
\]
从而有
\[
  \lim_{x\to a}\cos x=\cos a.
\]
同理可证
\[
  \lim_{x\to a}\sin x=\sin a.
\]
利用前面所证不等式，我们有
\[
  \cos x<\frac{\sin x}{x}<1,\qquad \forall x\in\left(0,\frac{\pi}{2}\right).
\]
利用偶函数性质，这一不等式在 $-\dfrac{\pi}{2}<x<0$ 时也成立. 再注意到，前面我们已经证明了 $\lim\limits_{x\to0}\cos x=1$，由夹逼收敛原理知 $\lim\limits_{x\to0}\dfrac{\sin x}{x}=1$ 成立。

\noindent\textbf{例 3.1.14}\quad 求极限
\[
  \lim_{x\to0}\frac{\sin(\sin x)}{x}.
\]

\noindent\textbf{解}\quad 由于 $|\sin x|<|x|$，所以 $\sin x\to0\ (x\to0)$. 这样，我们有
\[
\begin{aligned}
  \lim_{x\to0}\frac{\sin(\sin x)}{x}
  &=\lim_{x\to0}\frac{\sin(\sin x)}{\sin x}\frac{\sin x}{x}\\
  &=\lim_{x\to0}\frac{\sin(\sin x)}{\sin x}\lim_{x\to0}\frac{\sin x}{x}\\
  &=\lim_{t\to0}\frac{\sin t}{t}\lim_{x\to0}\frac{\sin x}{x}=1\qquad (t=\sin x).
\end{aligned}
\]

\noindent\textbf{第二个重要极限}\quad
\[
  \lim_{x\to\infty}\left(1+\frac1x\right)^x=e.
\]

\noindent\textbf{证明}\quad 先考虑 $x\to+\infty$ 的情形. 由于当 $x>1$ 时，有
\[
  1+\frac1{[x]+1}\le 1+\frac1x\le 1+\frac1{[x]},
\]
所以，有
\[
  \left(1+\frac1{[x]+1}\right)^x
  \le \left(1+\frac1x\right)^x
  \le \left(1+\frac1{[x]}\right)^x,
\]
从而有
\[
  \left(1+\frac1{[x]+1}\right)^{[x]}
  \le \left(1+\frac1x\right)^x
  \le \left(1+\frac1{[x]}\right)^{[x]+1}.
\]
这样，应用序列极限所得结果和夹逼收敛原理，知
\[
  \lim_{x\to+\infty}\left(1+\frac1x\right)^x=e.
\]

再考虑 $x\to-\infty$ 的情形. 令 $x=-y$，则 $y\to+\infty$，从而有
\[
\begin{aligned}
  \lim_{x\to-\infty}\left(1+\frac1x\right)^x
  &=\lim_{y\to+\infty}\left(1-\frac1y\right)^{-y}\\
  &=\lim_{y\to+\infty}\left(1+\frac1{y-1}\right)^{y-1}\left(1+\frac1{y-1}\right)=e.
\end{aligned}
\]
由极限与单侧极限的关系，得
\[
  \lim_{x\to\infty}\left(1+\frac1x\right)^x=e.
\]

\noindent\textbf{例 3.1.15}\quad 求极限
\[
  \lim_{x\to0}(1-2x)^{\frac1x}.
\]

\noindent\textbf{解}\quad 令 $x=\dfrac1t$，则 $x\to0$ 蕴涵着 $t\to\infty$. 于是
\[
\begin{aligned}
  \lim_{x\to0}(1-2x)^{\frac1x}
  &=\lim_{t\to\infty}\left(1-\frac2t\right)^t
  =\lim_{t\to\infty}\left(\frac{t-2}{t}\right)^t\\
  &=\lim_{t\to\infty}\left(1+\frac2{t-2}\right)^{-t}
  =\lim_{s\to\infty}\left(1+\frac1s\right)^{-(2s+2)}\\
  &=\frac1{\left(\lim\limits_{s\to\infty}\left(1+\dfrac1s\right)^s\right)^2}
  \lim_{s\to\infty}\left(1+\frac1s\right)^{-2}=e^{-2},
\end{aligned}
\]
这里 $s=\dfrac{t-2}{2}$。

\noindent\textbf{3. 函数的上、下极限}

作为本节的结束，我们简要地介绍一下函数的上、下极限. 为了避免繁琐，我们这里仅考虑当 $x\to x_0$ 时的情形，其他情形请读者自己补出。

设函数 $y=f(x)$ 在 $U_0(x_0,\delta_0)$ 内有定义. 对任意的 $0<\delta<\delta_0$，设
\[
  \ell(\delta)=\inf\{f(x):x\in U_0(x_0,\delta)\},
\]
\[
  h(\delta)=\sup\{f(x):x\in U_0(x_0,\delta)\},
\]
则当 $x\in U_0(x_0,\delta)$ 时，有 $\ell(\delta)\le f(x)\le h(\delta)$. 容易看出，作为 $(0,\delta_0)$ 上的函数，$\ell(\delta)$ 关于 $\delta$ 单调递减，而 $h(\delta)$ 关于 $\delta$ 单调递增. 因此它们的广义极限都存在. 令
\[
  \ell=\lim_{\delta\to0}\ell(\delta),\qquad h=\lim_{\delta\to0}h(\delta),
\]
则分别称之为当 $x$ 趋于 $x_0$ 时 $f(x)$ 的下极限和上极限，记为
\[
  \ell=\varliminf_{x\to x_0}f(x),\qquad h=\varlimsup_{x\to x_0}f(x).
\]

函数的上、下极限与序列的上、下极限具有相似的性质，请读者作为练习自己补出. 特别地，我们有

\noindent\textbf{定理 3.1.10}\quad 设函数 $f(x)$ 在 $U_0(x_0,\delta_0)(\delta_0>0)$ 内有定义，则 $\lim\limits_{x\to x_0}f(x)$ 存在的充分必要条件是
\[
  \varliminf_{x\to x_0}f(x)=\varlimsup_{x\to x_0}f(x).
\]

\section{函数的连续与间断}

\subsection{函数的连续与间断}

所谓函数在某点的连续与间断，是对该点函数值与其附近点的函数值的变化趋势是否衔接的一种刻画，是函数的一种局部属性，其定义如下：

\noindent\textbf{定义 3.2.1}\quad 设函数 $f(x)$ 在 $U(x_0,\delta_0)(\delta_0>0)$ 内有定义. 若有 $\lim\limits_{x\to x_0}f(x)=f(x_0)$，则称 $f(x)$ 在点 $x_0$ 连续，并称 $x_0$ 为 $f(x)$ 的一个连续点；否则称 $f(x)$ 在点 $x_0$ 间断（或不连续），并称 $x_0$ 为 $f(x)$ 的一个间断点（或不连续点）。

我们也可以用 $\varepsilon$-$\delta$ 语言来表述 $f(x)$ 在点 $x_0$ 处连续：若 $\forall\varepsilon>0$，$\exists\delta>0$，当 $x\in U(x_0,\delta)$ 时，有 $|f(x)-f(x_0)|<\varepsilon$，则称 $f(x)$ 在点 $x_0$ 连续。

设函数 $f(x)$ 在点 $x_0$ 处连续，则有
\[
  \lim_{x\to x_0}f(x)=f(x_0)=f\left(\lim_{x\to x_0}x\right).
\]
因此，函数在一点处连续可以看成在该点附近求函数值与求极限的顺序可以互换。

\noindent\textbf{例 3.2.1}\quad 在上一节例 3.1.2 之后，我们曾经指出，对任意的实数 $x_0$ 和任意的正整数 $k$，有 $\lim\limits_{x\to x_0}x^k=x_0^k$. 由函数极限的性质知，对任意的多项式函数
\[
  P(x)=a_nx^n+a_{n-1}x^{n-1}+\cdots+a_1x+a_0,
\]
有 $\lim\limits_{x\to x_0}P(x)=P(x_0)$. 这表明，多项式函数在 $\mathbb{R}$ 上的任意点处都连续。

\noindent\textbf{例 3.2.2}\quad 在上一节我们已经证明了
\[
  \lim_{x\to x_0}\sin x=\sin x_0,\qquad \lim_{x\to x_0}\cos x=\cos x_0.
\]
这里 $x_0$ 是任意给定的实数. 这表明，三角函数 $\sin x$ 和 $\cos x$ 在 $\mathbb{R}$ 上的任意点处都连续。

为了能够对闭区间上定义的函数讨论其端点处的连续性问题，我们引入：

\noindent\textbf{定义 3.2.2}\quad 若函数 $f(x)$ 在 $U^+(x_0,\delta_0)$ 上有定义，且 $f(x_0+0)=f(x_0)$，则称 $f(x)$ 在点 $x_0$ 右连续；若 $f(x)$ 在 $U^-(x_0,\delta_0)$ 上有定义，且 $f(x_0-0)=f(x_0)$，则称 $f(x)$ 在点 $x_0$ 左连续。

显然，$f(x)$ 在点 $x_0$ 处连续的充分必要条件是它在该点左、右连续。

\noindent\textbf{定义 3.2.3}\quad 设函数 $f(x)$ 在 $[a,b]$ 上有定义. 若对 $x\in(a,b)$，$f(x)$ 在点 $x$ 处连续，则称 $f(x)$ 在 $(a,b)$ 内连续，此时记为 $f(x)\in C(a,b)$；若 $f(x)\in C(a,b)$，而且在左端点 $a$ 右连续、在右端点 $b$ 左连续，则称 $f(x)$ 在 $[a,b]$ 上连续，此时记为 $f(x)\in C[a,b]$。

\noindent\textbf{注 1}\quad 设 $f(x)\in C(a,b)$，若 $\lim\limits_{x\to a+0}f(x)=A$ 及 $\lim\limits_{x\to b-0}f(x)=B$ 存在，则 $f(x)$ 可连续延拓到 $[a,b]$，即以下函数
\[
  \tilde f(x)=
  \begin{cases}
    A, & x=a,\\
    f(x), & x\in(a,b),\\
    B, & x=b
  \end{cases}
\]
是 $[a,b]$ 上的连续函数。

\noindent\textbf{注 2}\quad 若 $f(x)\in C[a,b]$，则 $f(x)$ 可连续延拓到 $(-\infty,+\infty)$. 事实上，以下函数就是 $f(x)$ 的一个连续延拓：
\[
  \tilde f(x)=
  \begin{cases}
    f(a), & x<a,\\
    f(x), & a\le x\le b,\\
    f(b), & x>b.
  \end{cases}
\]
请读者自己给出 $\tilde f(x)$ 的连续性证明。

如果 $f(x)$ 在 $x_0$ 处是间断的，则情况比较复杂. 一般地，我们可以将间断点分为以下两类：

(1) 若 $f(x_0+0)$ 与 $f(x_0-0)$ 都存在，则称 $x_0$ 为 $f(x)$ 的第一类间断点. 此时，若 $f(x_0+0)=f(x_0-0)\ne f(x_0)$，则称此间断点为可去间断点；否则称其为跳跃间断点。

(2) 若 $f(x_0+0)$ 与 $f(x_0-0)$ 至少有一个不存在时，则称 $x_0$ 为 $f(x)$ 的第二类间断点。

\noindent\textbf{例 3.2.3}\quad 考查函数
\[
  f(x)=
  \begin{cases}
    x\sin\dfrac1x, & x\ne0,\\
    1, & x=0
  \end{cases}
\]
在 $x=0$ 处的连续性。

\noindent\textbf{解}\quad 由 $\left|x\sin\dfrac1x\right|\le |x|$ 知，$\lim\limits_{x\to0}f(x)=0\ne f(0)$. 因此，$x=0$ 是 $f(x)$ 的可去间断点（参见图 3.2.1）。

\noindent\textbf{例 3.2.4}\quad 考查函数 $f(x)=\dfrac1x-\left[\dfrac1x\right]$ 在 $x=1$ 处的连续性。

\noindent\textbf{解}\quad 当 $x>1$ 时，有 $0<\dfrac1x<1$，于是，此时 $f(x)=\dfrac1x$，从而 $f(1+0)=1$；而当 $\dfrac12<x\le1$ 时，有 $1\le\dfrac1x<2$，于是，此时 $f(x)=\dfrac1x-1$，从而 $f(1-0)=0=f(1)$. 因此，$f(x)$ 在 $x=1$ 处左连续，右不连续，从而 $x=1$ 为跳跃间断点（参见图 3.2.2）。

\noindent\textbf{例 3.2.5}\quad 考查函数
\[
  f(x)=
  \begin{cases}
    e^{\frac1x}, & x\ne0,\\
    0, & x=0
  \end{cases}
\]
在 $x=0$ 处的连续性。

\noindent\textbf{解}\quad 由于
\[
  f(0+0)=+\infty,
\]
\[
  f(0-0)=0=f(0),
\]
所以 $x=0$ 是 $f(x)$ 的第二类间断点（参见图 3.2.3）。

\noindent\textbf{例 3.2.6}\quad 讨论函数
\[
  f(x)=
  \begin{cases}
    \sin\dfrac1x, & x\ne0,\\
    0, & x=0
  \end{cases}
\]
在 $x=0$ 处的连续性。

\noindent\textbf{解}\quad 在上一节我们已经证明了 $\lim\limits_{x\to0+0}f(x)$ 与 $\lim\limits_{x\to0-0}f(x)$ 都不存在. 因此，此时 $x=0$ 为 $f(x)$ 的第二类间断点（参见图 3.2.4）。

\noindent\textbf{例 3.2.7}\quad 讨论黎曼函数 $R(x)$（参考例 1.2.4）在 $(0,1)$ 内的连续性。

\noindent\textbf{解}\quad 任意取定 $x_0\in(0,1)$，我们来考查 $R(x)$ 在 $x_0$ 处的连续性。

$\forall\varepsilon>0$，取定正整数 $p$，使得 $1/p<\varepsilon$. 由于在 $(0,1)$ 内分母小于或等于 $p$ 的分数仅为有限个，在这有限个数中，必存在 $r\ne x_0$，使得 $r$ 到 $x_0$ 的距离是在所有和 $x_0$ 不相等的这些数中与 $x_0$ 的距离达到最小. 如果我们令 $\delta=\min\{|x_0-r|,x_0,1-x_0\}$，则对 $\forall x\in U_0(x_0,\delta)$，有
\[
  |R(x)-0|<1/p<\varepsilon,
\]
这表明 $\lim\limits_{x\to x_0}R(x)=0$. 由此可知：当 $x_0$ 为有理数时，$R(x)$ 在 $x_0$ 处不连续，是可去间断点；当 $x_0$ 为无理数时，$R(x)$ 在 $x_0$ 处连续。

\noindent\textbf{例 3.2.8}\quad 考查有理函数
\[
  f(x)=\frac{p(x)}{q(x)}
  =
  \frac{a_nx^n+a_{n-1}x^{n-1}+\cdots+a_1x+a_0}
       {b_mx^m+b_{m-1}x^{m-1}+\cdots+b_1x+b_0}
\]
的连续性以及 $x\to\infty$ 时的极限，其中 $a_nb_m\ne0$，而且 $p(x),q(x)$ 互素。

\noindent\textbf{解}\quad 设 $x_0\in(-\infty,+\infty)$. 当 $q(x_0)\ne0$ 时，有
\[
  \lim_{x\to x_0}\frac{p(x)}{q(x)}=\frac{p(x_0)}{q(x_0)},
\]
从而 $f(x)$ 在 $x_0$ 处连续；当 $q(x_0)=0$ 时，$f(x)$ 在 $x_0$ 处没有定义. 此外，由于 $p(x),q(x)$ 互素，所以 $p(x_0)\ne0$，从而
\[
  \lim_{x\to x_0}\frac{p(x)}{q(x)}=\infty.
\]
现考虑 $x\to\infty$ 的情况. 此时，我们有
\[
  \lim_{x\to\infty}f(x)
  =
  \lim_{x\to\infty}
  \left(
    x^{n-m}\frac{a_n+\dfrac{a_{n-1}}x+\cdots+\dfrac{a_0}{x^n}}
    {b_m+\dfrac{b_{m-1}}x+\cdots+\dfrac{b_0}{x^m}}
  \right)
  =
  \begin{cases}
    0, & n<m,\\[4pt]
    \dfrac{a_n}{b_m}, & n=m,\\[6pt]
    \infty, & n>m.
  \end{cases}
\]

\noindent\textbf{例 3.2.9}\quad 设函数 $f(x)$ 在 $(a,b)$ 上单调，则 $f(x)$ 在 $(a,b)$ 内只可能有第一类间断点。

\noindent\textbf{证明}\quad 不妨设 $f(x)$ 在 $(a,b)$ 上单调递减. 对 $\forall\xi\in(a,b)$，在 $(a,\xi)$ 中，$f(x)$ 单调递减且有下界 $f(\xi)$，因此 $\lim\limits_{x\to\xi-0}f(x)$ 存在；而在 $(\xi,b)$ 中，$f(x)$ 单调递减且有上界 $f(\xi)$，因此 $\lim\limits_{x\to\xi+0}f(x)$ 存在. 这样，若 $f(x)$ 在点 $\xi$ 处不连续，则 $\xi$ 只能是第一类间断点。

\noindent\textbf{注}\quad 由此例可以看出：若 $f(x)$ 在 $(a,b)$ 单调，则 $f(x)$ 在每一个间断点的取值具有跳跃性. 例如，假定 $f(x)$ 在 $(a,b)$ 单调递增，并且假设 $\xi\in(a,b)$ 为 $f(x)$ 的一个间断点，则必有 $f(\xi-0)<f(\xi+0)$. 因此 $f(x)$ 在 $(a,b)$ 中取不到 $(f(\xi-0),f(\xi+0))\setminus\{f(\xi)\}$ 中的值. 换句话说，$f(x)$ 在点 $\xi$ 处有一跳跃. 作为练习，请读者试构造一个单调函数，使其具有无穷多个间断点。

\subsection{连续函数的性质}

首先，由函数在一点连续的定义和函数极限的性质，读者不难得到：

(1) 连续函数是局部有界的，即若 $f(x)$ 在点 $x_0$ 处连续，则必存在 $\delta>0$，使得 $f(x)$ 在 $U(x_0,\delta)$ 上有界。

(2) 连续函数具有局部保号性，即：若 $f(x)$ 在点 $x_0$ 处连续，而且 $f(x_0)>0$，则必存在 $\delta>0$，使得 $f(x)>0$ 对一切 $x\in U(x_0,\delta)$ 成立. 以后我们要经常用到的结论是：对任意的 $0<A<f(x_0)$（如 $A=f(x_0)/2$），总存在 $\delta_0>0$，使得当 $x\in U(x_0,\delta_0)$ 时，有 $f(x)>A$。

(3) 连续函数经四则运算后仍然连续，即：若 $f(x)$ 和 $g(x)$ 在点 $x_0$ 处连续，则函数 $f(x)\pm g(x)$，$f(x)g(x)$ 和 $\dfrac{f(x)}{g(x)}$（$g(x_0)\ne0$）仍然在点 $x_0$ 处连续。

此外，连续函数的复合函数和反函数也是连续的。

\noindent\textbf{定理 3.2.1（复合函数的连续性）}\quad 设 $u=g(x)$ 在点 $x_0$ 连续，$y=f(u)$ 在点 $u_0=g(x_0)$ 连续，则复合函数 $f(g(x))$ 在点 $x_0$ 连续。

\noindent\textbf{证明}\quad $\forall\varepsilon>0$，由 $f(u)$ 在点 $u_0$ 连续知，$\exists\eta>0$，当 $|u-u_0|<\eta$ 时，有
\[
  |f(u)-f(u_0)|<\varepsilon.
\]
对 $\eta>0$，再由 $u=g(x)$ 在点 $x_0$ 连续知，$\exists\delta>0$，当 $|x-x_0|<\delta$ 时，有
\[
  |u-u_0|=|g(x)-g(x_0)|<\eta,
\]
从而当 $|x-x_0|<\delta$ 时，有
\[
  |f(g(x))-f(g(x_0))|<\varepsilon.
\]
因此，$f(g(x))$ 在点 $x_0$ 连续。

\noindent\textbf{定理 3.2.2（反函数的连续性）}\quad 设 $f(x)$ 是区间 $I$ 上严格单调的连续函数，则其反函数 $x=f^{-1}(y)$ 在 $f(I)$ 上连续。

\noindent\textbf{证明}\quad 由 $f(x)$ 的严格单调性，知 $f^{-1}(y)$ 存在. 现不妨设 $f(x)$ 严格递增. 显然，当 $f(x)$ 严格递增时，$f^{-1}(x)$ 必定严格递增. 另外，由 $f(x)$ 的连续性，我们可以推出 $f(I)$ 必定是一个区间（在下节中，我们将证明任何区间上连续函数的值域都是一个区间）. 事实上，倘若 $f(I)$ 不是一个区间，则由 $f(x)$ 的连续性，必存在一个开区间 $(\alpha,\beta)$ 使得 $f(I)\cap(\alpha,\beta)=\varnothing$，并且 $\alpha,\beta\in f(I)$. 设 $f(x_0)=\alpha$. 由 $f(x)$ 的单调性有
\[
  f(x_0-0)\le f(x_0)<\beta\le f(x_0+0).
\]
这与 $f(x)$ 在 $x_0$ 的连续性矛盾。

现任取 $y_0\in f(I)$，则存在 $x_0\in I$，使得 $x_0=f^{-1}(y_0)$. 不妨假定 $x_0$ 不是 $I$ 的端点，否则考查左、右连续即可。

$\forall\varepsilon>0$，存在 $x_1,x_2\in I$，使得 $x_0-\varepsilon<x_1<x_0<x_2<x_0+\varepsilon$. 记 $y_1=f(x_1),y_2=f(x_2)$，则有
\[
  y_1=f(x_1)<f(x_0)=y_0<f(x_2)=y_2.
\]
令 $\delta=\min\{y_0-y_1,y_2-y_0\}>0$，则当 $y\in(y_0-\delta,y_0+\delta)$ 时，有 $y_1<y<y_2$，从而由 $f(I)$ 是一个区间推知，$\exists x\in I$，使得 $y=f(x)$，即 $y\in f(I)$。
于是
\[
  x_0-\varepsilon<x_1=f^{-1}(y_1)<f^{-1}(y)<f^{-1}(y_2)=x_2<x_0+\varepsilon
\]
即
\[
  |f^{-1}(y)-f^{-1}(y_0)|=|f^{-1}(y)-x_0|<\varepsilon.
\]
这就证明了 $f^{-1}(y)$ 在点 $y_0$ 处连续. 注意到 $y_0\in f(I)$ 的任意性，即知 $f^{-1}(y)$ 在 $f(I)$ 上连续。

\subsection{初等函数的连续性}

首先我们来证明，对 $a>1$，指数函数 $f(x)=a^x$ 在 $(-\infty,+\infty)$ 上连续. 我们回忆一下，实际上，在中学数学的学习中，我们并不知道当 $x$ 是无理数时 $a^x$ 是如何定义的. 下面我们给出它的严格定义：

对任何正有理数 $x=\dfrac qp$（其中 $p,q$ 是互素的正整数），定义 $a^x=(\sqrt[p]{a})^q$；对于任何负有理数 $x$，定义 $a^x=\dfrac1{a^{-x}}$，并定义 $a^0=1$. 容易证明，$\forall r,s\in\mathbb Q$，有
\[
  \text{(1) } a^ra^s=a^{r+s};
\]
\[
  \text{(2) 若 } r<s,\text{ 则 } a^r<a^s.
\]
而当 $x$ 为一无理数时，定义
\[
  a^x=\sup\{a^r:r\text{ 为小于 }x\text{ 的有理数}\}.
\]
这样，对一切 $x\in(-\infty,+\infty)$，$a^x$ 均有定义。

现在我们来证明 $a^x$ 在 $(-\infty,+\infty)$ 上严格递增. 设 $x_1<x_2$，则存在有理数 $q_1,q_2$，使得
\[
  x_1\le q_1<q_2\le x_2,
\]
从而
\[
  a^{x_1}=\sup\{a^q:q\text{ 为小于 }x_1\text{ 的有理数}\}
  \le a^{q_1}<a^{q_2}
  \le \sup\{a^q:q\text{ 为小于 }x_2\text{ 的有理数}\}=a^{x_2}.
\]

再来证明 $a^x\in C(-\infty,+\infty)$. $\forall x_0\in(-\infty,+\infty)$，$\forall\varepsilon>0$，先取正整数 $N$，使得 $a<(1+\varepsilon)^{N/2}$；再取正整数 $q$，使得 $q\le Nx_0<q+1$；然后令
\[
  q_1=\frac{q-1}{N},\quad q_2=\frac{q+1}{N}.
\]
这样，由 $a^x$ 的严格递增性知
\[
  a^{q_1}\le f(x_0-0)\le f(x_0)\le f(x_0+0)\le a^{q_2},
\]
从而有
\[
  1\le \frac{f(x_0+0)}{f(x_0-0)}
  \le \frac{a^{q_2}}{a^{q_1}}
  =a^{q_2-q_1}=a^{2/N}<1+\varepsilon.
\]
由 $\varepsilon>0$ 的任意性，即知
\[
  f(x_0+0)=f(x_0-0)=f(x_0)=a^{x_0}.
\]
这样，我们就证明了 $a^x\in C(-\infty,+\infty)$。

对于 $0<a<1$，我们定义 $a^x=\dfrac1{(1/a)^x}$，则由连续函数的性质知，此时亦有 $a^x\in C(-\infty,+\infty)$；至于 $a=1$ 的情形，我们规定 $a^x=1$. 这样一来，对 $\forall a\in(0,+\infty)$，$y=a^x$ 均有定义，而且 $y=a^x\in C(-\infty,+\infty)$。

由反函数的连续性知，对任意的 $a>0$，对数函数 $y=\log_a x\in C(0,+\infty)$；再由复合函数连续性知，幂函数 $y=x^\alpha=e^{\alpha\ln x}\in C(0,+\infty)$。

上一节我们已经证明了，对任意的实数 $a$，有 $\lim\limits_{x\to a}\cos x=\cos a$，从而 $\cos x\in C(-\infty,+\infty)$. 同理可证 $\sin x\in C(-\infty,+\infty)$. 再利用连续函数的性质知，所有的三角函数及反三角函数在其定义域内都是连续的。

这样一来，我们已经证明了所有的基本初等函数在其定义域内都是连续的. 由于初等函数是由基本初等函数经过有限次四则运算和复合运算所得，利用连续函数的性质，我们有

\noindent\textbf{定理 3.2.3}\quad 初等函数在其定义域内是连续的。

\noindent\textbf{例 3.2.10}\quad 设 $\lim\limits_{x\to x_0}u(x)=a>0$，$\lim\limits_{x\to x_0}v(x)=b$，则
\[
  \lim_{x\to x_0}u(x)^{v(x)}=a^b.
\]

\noindent\textbf{证明}\quad 定义 $u(x_0)=a$，$v(x_0)=b$，则 $u(x),v(x)$ 在点 $x_0$ 处都连续，从而
\[
  \lim_{x\to x_0}u(x)^{v(x)}
  =
  \lim_{x\to x_0}e^{v(x)\ln u(x)}
  =
  e^{v(x_0)\ln u(x_0)}
  =
  u(x_0)^{v(x_0)}
  =
  a^b.
\]

\noindent\textbf{例 3.2.11}\quad 求极限
\[
  \lim_{x\to+\infty}\left(1+\sin\frac1x\right)^x.
\]

\noindent\textbf{解}\quad 由
\[
  \left(1+\sin\frac1x\right)^x
  =
  \left(1+\sin\frac1x\right)^{\frac1{\sin(1/x)}\cdot\frac{\sin(1/x)}{1/x}},
\]
\[
  \lim_{x\to+\infty}\left(1+\sin\frac1x\right)^{\frac1{\sin(1/x)}}=e,
  \qquad
  \lim_{x\to+\infty}\frac{\sin(1/x)}{1/x}=1,
\]
得
\[
  \lim_{x\to+\infty}\left(1+\sin\frac1x\right)^x=e^1=e.
\]

\section{闭区间上连续函数的基本性质}

上一节中，我们讨论了函数 $f(x)$ 在其连续点 $x_0$ 邻近的局部性质. 例如，若一个函数 $f(x)$ 在点 $x_0$ 连续且满足 $f(x_0)<\eta$，则我们能断定在点 $x_0$ 邻近有 $f(x)<\eta$，即存在 $\delta>0$，使得 $f(x)<\eta$ 对一切的 $x\in(x_0-\delta,x_0+\delta)$ 成立. 如果一个函数 $f(x)\in C[a,b]$，那么 $f(x)$ 在整个闭区间上能具有什么样的整体性质呢？本节将讨论这一重要问题。

\noindent\textbf{定理 3.3.1（有界性）}\quad 设函数 $f(x)\in C[a,b]$，则 $f(x)$ 在 $[a,b]$ 上有界。

\noindent\textbf{证明}\quad 倘若 $f(x)$ 在 $[a,b]$ 上无界，则将 $[a,b]$ 等分成两个闭区间，$f(x)$ 必在其中一个子区间上无界，取定这样的一个区间并将其设为 $[a_1,b_1]$. 同理，再将 $[a_1,b_1]$ 二等分，$f(x)$ 必在其中一个子区间上无界，取定这样的一个区间并将其记为 $[a_2,b_2]$. 如此进行下去，我们就得到一个闭区间列 $\{[a_n,b_n]\}$，满足：
\[
  \text{(1) } [a_{n+1},b_{n+1}]\subset [a_n,b_n];
\]
\[
  \text{(2) } b_{n+1}-a_{n+1}=\frac12(b_n-a_n);
\]
\[
  \text{(3) } f(x)\text{ 在每个 }[a_n,b_n]\text{ 上都无界.}
\]
于是，由闭区间套定理知，存在唯一的 $\xi\in\bigcap\limits_{n=1}^{+\infty}[a_n,b_n]$. 再由 $f(x)\in C[a,b]$ 和 $\xi\in[a_1,b_1]\subset[a,b]$ 知，存在 $\delta>0$，使得 $f(x)$ 在 $U(\xi,\delta)\cap[a,b]$ 上有界；而条件 (2) 又蕴涵着，当 $n$ 充分大时，有 $[a_n,b_n]\subset U(\xi,\delta)$，从而 $f(x)$ 在 $[a_n,b_n]$ 上有界. 这与条件 (3) 矛盾。

\noindent\textbf{例 3.3.1}\quad 函数 $y=\dfrac1x$ 在开区间 $(0,1)$ 上连续，但它在此区间上无界。

\noindent\textbf{思考题}\quad 若不假定 $f(x)$ 在 $[a,b]$ 上连续，仅仅假定 $\forall x_0\in[a,b]$，$\lim\limits_{x\to x_0}f(x)$ 存在（在端点处单侧极限存在），能否推出 $f(x)$ 在 $[a,b]$ 上有界。

设 $f(x)\in C[a,b]$. 由于 $f(x)$ 在 $[a,b]$ 上有界，故 $\sup\limits_{x\in[a,b]}\{f(x)\}$ 与 $\inf\limits_{x\in[a,b]}\{f(x)\}$ 都为有限数. 下面的定理告诉我们它们将属于 $f(x)$ 的值域。

\noindent\textbf{定理 3.3.2（最值定理）}\quad 设 $f(x)\in C[a,b]$，则 $f(x)$ 在 $[a,b]$ 上必有最小值和最大值，即存在 $\xi,\zeta\in[a,b]$，使得 $f(\xi)\le f(x)\le f(\zeta)$ 对一切的 $x\in[a,b]$ 成立。

\noindent\textbf{证明}\quad 令 $M=\sup\limits_{x\in[a,b]}\{f(x)\}$，则由上确界的定义知，存在一个序列 $\{x_n\}\subset[a,b]$，使得 $f(x_n)\to M$（$n\to\infty$）. 由于 $\{x_n\}$ 是有界序列，所以必有一收敛子列 $\{x_{n_k}\}$. 设 $\lim\limits_{k\to\infty}x_{n_k}=\zeta$，则 $\zeta\in[a,b]$. 再由 $f(x)\in C[a,b]$ 知
\[
  f(\zeta)=\lim_{k\to\infty}f(x_{n_k})=M.
\]
同理可证存在 $\xi\in[a,b]$，使得
\[
  f(\xi)=m=\inf_{x\in[a,b]}\{f(x)\}.
\]
此定理表明：若 $f(x)\in C[a,b]$，则存在 $\xi,\zeta\in[a,b]$，使得
\[
  f(\xi)=\min_{x\in[a,b]}\{f(x)\},\qquad
  f(\zeta)=\max_{x\in[a,b]}\{f(x)\}.
\]

\noindent\textbf{例 3.3.2}\quad 设 $f(x)\in C[a,b]$，且 $\forall x\in[a,b]$，有 $f(x)>0$. 证明存在正数 $\eta$，使得 $f(x)\ge\eta$ 对一切的 $x\in[a,b]$ 成立。

\noindent\textbf{证明}\quad 由 $f(x)\in C[a,b]$，根据定理 3.3.2 知，$\exists\xi\in[a,b]$，使
\[
  f(\xi)=\min_{x\in[a,b]}\{f(x)\}>0,
\]
从而取 $\eta=f(\xi)$ 即可。

\noindent\textbf{定理 3.3.3（介值定理）}\quad 设 $f(x)\in C[a,b]$，记 $m=\min\limits_{x\in[a,b]}\{f(x)\}$，$M=\max\limits_{x\in[a,b]}\{f(x)\}$，则 $f([a,b])=[m,M]$，即对 $\forall\eta\in(m,M)$，$\exists\xi\in[a,b]$，使得 $f(\xi)=\eta$。

\noindent\textbf{证明}\quad 由定理 3.3.2 知，$\exists x_1,x_2\in[a,b]$，使得 $f(x_1)=m$，$f(x_2)=M$. 下面证明对 $\forall\eta\in(m,M)$，$\exists\xi\in[a,b]$，使得 $f(\xi)=\eta$。

不妨设 $x_1<x_2$，并定义
\[
  E=\{x\in[x_1,x_2]:f(x)>\eta\}.
\]
显然有 $x_1\notin E$，$x_2\in E$. 现令 $E$ 的下确界为 $\xi$，则有 $f(\xi)=\eta$. 事实上，若不然，则 $f(\xi)>\eta$. 由连续函数的局部保号性知，存在 $\xi$ 的一个邻域 $U(\xi,\delta)$，使得 $f(x)>\eta$ 对一切的 $x\in U(\xi,\delta)$ 成立. 由于 $\xi>x_1$，所以 $[x_1,x_2]\cap(\xi-\delta,\xi)\ne\varnothing$，现取 $\xi'\in[x_1,x_2]\cap(\xi-\delta,\xi)$，则有 $f(\xi')>\eta$，从而有 $\xi'<\xi$ 且 $\xi'\in E$. 这与 $\xi$ 是 $E$ 的下确界矛盾。

\noindent\textbf{例 3.3.3}\quad 设 $f(x)\in C[a,b]$，$x_j\in[a,b]$，$j=1,2,\cdots,n$. 证明 $\exists\xi\in[a,b]$，使得
\[
  f(\xi)=\frac1n\sum_{j=1}^{n}f(x_j).
\]

\noindent\textbf{证明}\quad 记
\[
  m=\min_{x\in[a,b]}\{f(x)\},\qquad
  M=\max_{x\in[a,b]}\{f(x)\},
\]
则有
\[
  m\le\frac1n\sum_{j=1}^{n}f(x_j)\le M.
\]
于是，由介值定理知，$\exists\xi\in[a,b]$，使得
\[
  f(\xi)=\frac1n\sum_{j=1}^{n}f(x_j).
\]

\noindent\textbf{注}\quad 将例 3.3.3 中的 $[a,b]$ 换成开区间 $(a,b)$，其结论仍真。

\noindent\textbf{定理 3.3.4（零点存在定理）}\quad 设 $f(x)$ 在区间 $I$ 上连续. 若 $\alpha,\beta\in I$，$\alpha<\beta$，满足 $f(\alpha)f(\beta)<0$，则存在 $\xi\in(\alpha,\beta)$，使得 $f(\xi)=0$。

\noindent\textbf{证明}\quad 不妨设 $f(\alpha)<0$，$f(\beta)>0$，则有
\[
  m=\min_{x\in[\alpha,\beta]}\{f(x)\}<0<\max_{x\in[\alpha,\beta]}\{f(x)\}=M.
\]
在闭区间 $[\alpha,\beta]$ 上应用介值定理，知存在 $\xi\in[\alpha,\beta]$，使得 $f(\xi)=0$. 但 $f(\alpha)<0$，而 $f(\beta)>0$，所以必有 $\xi\in(\alpha,\beta)$ 成立。

\noindent\textbf{注}\quad 这一定理表明：任何使连续函数 $f(x)$ 的函数值异号的两点之间必有方程 $f(x)=0$ 的一个根. 我们常可以用它来证明非线性方程之解的存在性，或者用它来求方程的近似解. 显然，定理 3.3.4 是定理 3.3.3 的直接推论. 但由于零点存在定理有特别的重要性，我们还是将它叙述成一个单独的定理。

\noindent\textbf{例 3.3.4}\quad 证明三次方程 $ax^3+bx^2+cx+d=0$ 至少有一实根。

\noindent\textbf{证明}\quad 无妨设 $a>0$. 由于
\[
  f(x)=ax^3+bx^2+cx+d
  =
  ax^3\left(1+\frac b{ax}+\frac c{ax^2}+\frac d{ax^3}\right),
\]
我们有 $\lim\limits_{x\to-\infty}f(x)=-\infty$ 和 $\lim\limits_{x\to+\infty}f(x)=+\infty$. 因此易知，必有 $\alpha<0$ 和 $\beta>0$，使 $f(\alpha)<0$，$f(\beta)>0$. 在闭区间 $[\alpha,\beta]$ 上应用定理 3.3.4 即知在 $(\alpha,\beta)$ 内必有 $f(x)$ 的一个零点，即该三次方程至少有一个实根。

\noindent\textbf{例 3.3.5}\quad 证明方程 $x^4+2x-1=0$ 在 $(0,1)$ 中必有一根，并求它的一个近似解。

\noindent\textbf{证明}\quad 令 $f(x)=x^4+2x-1$，则有 $f(0)=-1<0$，$f(1)=2>0$. 于是，在 $(0,1)$ 之内必有该方程的一个根. 取 $[0,1]$ 的中点 $\dfrac12$，计算可得 $f\left(\dfrac12\right)=0.0625>0$. 于是，在 $\left(0,\dfrac12\right)$ 之内必有该方程的一个根. 再取 $\left[0,\dfrac12\right]$ 的中点 $\dfrac14$，计算可得 $f\left(\dfrac14\right)=-0.4961<0$. 于是，在 $\left(\dfrac14,\dfrac12\right)$ 之内必有该方程的一个根. 如此进行，我们可以断定，在 $\left(\dfrac{121}{256},\dfrac{61}{128}\right)$ 之内必有该方程的一个根，因此我们可取其中点
\[
  \hat x=\frac{243}{512}\approx0.4746
\]
作为该方程的一个近似解，其误差为
\[
  |x-\hat x|<\frac1{512}\approx0.002.
\]

\noindent\textbf{例 3.3.6}\quad 设函数 $f(x)\in C[a,b]$，且 $f^{-1}(x)$ 存在. 证明 $f(x)$ 必在 $[a,b]$ 上严格单调。

\noindent\textbf{证明}\quad 倘若 $f(x)$ 在 $[a,b]$ 不严格单调，必存在
\[
  a\le x_1<x_2<x_3\le b,
\]
使得以下两种情形之一发生：
\[
  \text{(1) } f(x_1)<f(x_2)\text{ 且 }f(x_2)>f(x_3);
\]
\[
  \text{(2) } f(x_1)>f(x_2)\text{ 且 }f(x_2)<f(x_3).
\]
不妨设情形 (1) 发生，并且取 $\eta$ 满足
\[
  \max\{f(x_1),f(x_3)\}<\eta<f(x_2),
\]
则由介值定理推知，存在 $\xi_1\in(x_1,x_2)$ 和 $\xi_2\in(x_2,x_3)$，使得 $f(\xi_1)=\eta=f(\xi_2)$. 这与 $f^{-1}(x)$ 存在相矛盾。

\noindent\textbf{注}\quad 前面我们已经指出（见例 3.2.9 的注）：若单调函数 $f(x)$ 在 $I$ 内不连续，则它的取值不能是一个区间. 将这一事实与连续函数的介值性定理相结合便有：若 $f(x)$ 在区间 $I$ 上单调，则 $f(x)$ 在 $I$ 上连续的充分必要条件是 $f(I)$ 是一个区间。

以下我们来讨论函数的一致连续性问题. 为此先来考查函数 $f(x)=\dfrac1x$，$x\in(0,1)$. 我们知道 $f(x)$ 在 $(0,1)$ 是连续的. 换句话说，对每个 $x_0\in(0,1)$，对任意给定的 $\varepsilon>0$，都存在 $\delta>0$，使得当 $|x-x_0|<\delta$ 时，有 $\left|\dfrac1x-\dfrac1{x_0}\right|<\varepsilon$. 这里有个值得注意的现象是，对于同一个 $\varepsilon$，当 $x_0$ 与 $0$ 越近时，相应得到的 $\delta$ 也就越小. 换句话说，我们不可能找到一个仅依赖 $\varepsilon$，而不依赖 $x_0$ 的 $\delta$，使得对所有的 $x_0\in(0,1)$，有 $\left|\dfrac1x-\dfrac1{x_0}\right|<\varepsilon$。

基于这一现象，我们引入如下的概念：

\noindent\textbf{定义 3.3.1}\quad 设函数 $f(x)$ 在区间 $I$ 上有定义. 若 $\forall\varepsilon>0$，$\exists\delta>0$，当 $x_1,x_2\in I$ 且 $|x_1-x_2|<\delta$ 时，有 $|f(x_1)-f(x_2)|<\varepsilon$，则称 $f(x)$ 在 $I$ 上一致连续。

显然，$f(x)$ 在 $I$ 上一致连续，它必在 $I$ 上连续。

\noindent\textbf{例 3.3.7}\quad 证明函数 $y=\sin x$ 在 $(-\infty,+\infty)$ 上一致连续。

\noindent\textbf{证明}\quad 由于对任意的 $x_1,x_2\in(-\infty,+\infty)$ 有
\[
  |\sin x_1-\sin x_2|
  =
  2\left|\cos\frac{x_1+x_2}{2}\right|
  \left|\sin\frac{x_1-x_2}{2}\right|
  \le |x_1-x_2|,
\]
所以 $\forall\varepsilon>0$，取 $\delta=\varepsilon$，当 $x_1,x_2\in(-\infty,+\infty)$ 且 $|x_1-x_2|<\delta$ 时，有 $|\sin x_1-\sin x_2|<\varepsilon$. 这就证明了 $y=\sin x$ 在 $(-\infty,+\infty)$ 上一致连续。

对于判别一个函数是否一致连续，下面的定理有时是非常方便的。

\noindent\textbf{定理 3.3.5}\quad 设函数 $f(x)$ 在区间 $I$ 上有定义，则 $f(x)$ 在 $I$ 一致连续的充分必要条件是：对任意的两个序列 $\{x_n'\}\subset I,\{x_n''\}\subset I$，若它们满足 $\lim\limits_{n\to\infty}(x_n'-x_n'')=0$，必有
\[
  \lim_{n\to\infty}[f(x_n')-f(x_n'')]=0.
\]
此定理的证明留给读者自己完成。

\noindent\textbf{例 3.3.8}\quad 证明函数 $f(x)=\dfrac1x$ 在 $(0,+\infty)$ 上不一致连续；但对任意的 $x_0>0$，它在 $[x_0,+\infty)$ 上一致连续。

\noindent\textbf{证明}\quad 由定理 3.3.5 我们注意到，“$f(x)$ 在 $I$ 上不一致连续”与以下断言等价：存在 $\varepsilon_0>0$ 及序列 $\{x_n'\},\{x_n''\}\subset I$ 满足 $x_n'-x_n''\to0$（$n\to\infty$）及 $|f(x_n')-f(x_n'')|\ge\varepsilon_0$。

对 $f(x)=\dfrac1x$ 而言，我们取 $\varepsilon_0=1$，$x_n'=\dfrac1n$，$x_n''=\dfrac1{n+1}$（$n=1,2,\cdots$），则有
\[
  x_n'-x_n''\to0\quad(n\to\infty),
\]
且
\[
  \left|\frac1{x_n'}-\frac1{x_n''}\right|=1.
\]
此即证明了 $\dfrac1x$ 在 $(0,+\infty)$ 上不一致连续。

对任意给定的 $x_0>0$，由于对 $x_1,x_2\in[x_0,+\infty)$，有
\[
  \left|\frac1{x_1}-\frac1{x_2}\right|
  =
  \frac{|x_1-x_2|}{x_1x_2}
  \le \frac{|x_1-x_2|}{x_0^2},
\]
所以 $\forall\varepsilon>0$，取 $\delta=x_0^2\varepsilon>0$，则当 $x_1,x_2\in[x_0,+\infty)$ 且 $|x_1-x_2|<\delta$ 时，就有
\[
  \left|\frac1{x_1}-\frac1{x_2}\right|<\varepsilon.
\]
此时的 $\delta$ 只依赖于 $\varepsilon$，所以 $f(x)$ 在 $[x_0,+\infty)$ 上一致连续。

\noindent\textbf{定理 3.3.6（康托尔定理）}\quad 设函数 $f(x)\in C[a,b]$，则 $f(x)$ 在闭区间 $[a,b]$ 上一致连续。

\noindent\textbf{证明}\quad 定义
\[
  F(x)=
  \begin{cases}
    f(a), & x\in(-\infty,a),\\
    f(x), & x\in[a,b],\\
    f(b), & x\in(b,+\infty),
  \end{cases}
\]
则 $F(x)\in C(-\infty,+\infty)$. $\forall\varepsilon>0$ 和 $x\in[a,b]$，由 $F(x)$ 在点 $x$ 处的连续性知，$\exists\delta_x>0$，当 $x',x''\in U(x,2\delta_x)$ 时，有
\[
  |F(x')-F(x'')|<\varepsilon.
\]
显然，开区间簇 $\{U(x,\delta_x):x\in[a,b]\}$ 是闭区间 $[a,b]$ 的一个开覆盖. 由有限覆盖定理知，必有有限多个这样的开区间
\[
  U(x_1,\delta_{x_1}),\quad U(x_2,\delta_{x_2}),\quad \cdots,\quad U(x_m,\delta_{x_m}),
\]
它们也完全覆盖了 $[a,b]$. 令 $\delta=\min\{\delta_{x_j}:j=1,2,\cdots,m\}$，则 $\delta>0$，而且当 $x',x''\in[a,b]$ 且 $|x'-x''|<\delta$ 时，必存在某一 $U(x_j,\delta_{x_j})$，使得 $x'\in U(x_j,\delta_{x_j})$，即 $|x'-x_j|<\delta_{x_j}$，从而
\[
  |x''-x_j|\le |x''-x'|+|x'-x_j|<\delta+\delta_{x_j}\le2\delta_{x_j},
\]
于是有 $x',x''\in U(x_j,2\delta_{x_j})$. 这样，由 $\delta_{x_j}$ 的定义，有
\[
  |f(x')-f(x'')|=|F(x')-F(x'')|<\varepsilon.
\]
注意到 $\delta$ 只与 $\varepsilon$ 有关而与 $x',x''$ 无关，这就证明了 $f(x)$ 在 $[a,b]$ 上一致连续。

\noindent\textbf{例 3.3.9}\quad 设函数 $f(x)$ 在 $(a,b)$ 上连续. 证明：$f(x)$ 在 $(a,b)$ 上一致连续的充分必要条件是 $\lim\limits_{x\to a+0}f(x)$ 与 $\lim\limits_{x\to b-0}f(x)$ 都存在。

\noindent\textbf{证明}\quad 必要性\quad 若 $f(x)$ 在 $(a,b)$ 上一致连续，由一致连续的定义及函数极限的柯西收敛准则知，$\lim\limits_{x\to a+0}f(x)$ 与 $\lim\limits_{x\to b-0}f(x)$ 都存在。

充分性\quad 令
\[
  f(a)=\lim_{x\to a+0}f(x),\qquad f(b)=\lim_{x\to b-0}f(x),
\]
并且定义
\[
  \tilde f(x)=
  \begin{cases}
    f(a), & x=a,\\
    f(x), & x\in(a,b),\\
    f(b), & x=b,
  \end{cases}
\]
则 $\tilde f\in C[a,b]$. 于是由定理 3.3.6 知，它在 $[a,b]$ 上一致连续，从而 $f(x)$ 在 $(a,b)$ 上一致连续。

\noindent\textbf{注}\quad 从以上例子可知，要考查一个连续函数在有穷区间上是否一致连续，只需考查它在端点的单侧极限是否存在即可. 从这个意义上来讲，连续函数在有穷区间的一致连续性问题已经圆满解决. 特别地，若 $f(x)$ 在有穷区间 $I$ 上无界，则它在 $I$ 上必定不一致连续。

当函数 $f(x),g(x)$ 在区间 $I$ 上一致连续时，显然 $f(x)\pm g(x)$ 在区间 $I$ 上也一致连续。

\noindent\textbf{思考题}\quad 设 $f(x),g(x)$ 在区间 $I$ 上一致连续，是否 $f(x)g(x)$ 也在 $I$ 上一致连续？

\noindent\textbf{例 3.3.10}\quad 证明 $f(x)=\sqrt{x}$ 在 $[0,+\infty)$ 上一致连续。

\noindent\textbf{证明}\quad 显然 $f(x)$ 在 $[0,+\infty)$ 连续. 下面我们证明 $f(x)=\sqrt{x}$ 在 $[0,+\infty)$ 一致连续。

事实上，任取 $\varepsilon>0$，取 $\delta_1=\varepsilon$，则当 $x',x''\in[1,+\infty)$ 且 $|x'-x''|<\delta_1$ 时，有
\[
  |f(x')-f(x'')|
  =
  |\sqrt{x'}-\sqrt{x''}|
  =
  \frac{|x'-x''|}{|\sqrt{x'}+\sqrt{x''}|}
  \le |x'-x''|<\varepsilon.
\]
其次，由康托尔定理，$f(x)$ 在 $[0,2]$ 上连续，从而一致连续. 因此对上述 $\varepsilon>0$，$\exists\delta_2>0$，使得当 $x',x''\in[0,2]$ 且 $|x'-x''|<\delta_2$ 时，有
\[
  |f(x')-f(x'')|<\varepsilon.
\]
因此取 $\delta=\min\left\{\delta_1,\delta_2,\dfrac12\right\}$，对任意的 $x',x''\in[0,+\infty)$，当 $|x'-x''|<\delta$ 时，必有 $x',x''$ 同时落在 $[0,2]$ 或 $[1,+\infty)$ 上. 因此，由上所证，有
\[
  |f(x')-f(x'')|<\varepsilon.
\]
这就证明了 $f(x)=\sqrt{x}$ 在 $[0,+\infty)$ 上一致连续。

\section{无穷小量与无穷大量的阶}

首先，我们以 $x\to x_0$ 为规范，给出与无穷小量和无穷大量有关的一些定义和记号。

\noindent\textbf{定义 3.4.1}\quad 设函数 $f(x)$ 在 $U_0(x_0,\delta_0)(\delta_0>0)$ 上有定义. 若 $\lim\limits_{x\to x_0}f(x)=0$，则称 $f(x)$ 为 $x\to x_0$ 时的一个无穷小量；若 $\lim\limits_{x\to x_0}f(x)=\infty$，则称 $f(x)$ 为 $x\to x_0$ 时的一个无穷大量。

无穷小量有特殊的意义，在数学分析产生与发展的历史上，它起过特别重要的作用. 这是因为无穷小量不仅是一种特殊的有极限的变量，而且由于：$x\to x_0$ 时 $f(x)$ 以 $A$ 为极限的充分必要条件是 $f(x)-A$ 为 $x\to x_0$ 时的一个无穷小量，从而无穷小量可以表征一般有极限的变量。

关于无穷小量与无穷大量之间有如下的关系：

\noindent\textbf{定理 3.4.1}\quad 设函数 $f(x)$ 在 $U_0(x_0,\delta_0)(\delta_0>0)$ 上有定义且恒不为零，则 $f(x)$ 为 $x\to x_0$ 时的一个无穷小量的充分必要条件是 $\dfrac1{f(x)}$ 为 $x\to x_0$ 时的一个无穷大量。

\noindent\textbf{证明略.}

在前面我们已经知道 $\left\{\dfrac1n\right\}$ 和 $\left\{\dfrac1{n^n}\right\}$ 都是无穷小量，而且后者趋于零的速度比前者要“快”. 如何来刻画它们趋于零的速度是一个值得研究的问题. 为了便于比较无穷小量（无穷大量），我们引入：

\noindent\textbf{定义 3.4.2}\quad 设 $f(x)$ 和 $g(x)$ 都是当 $x\to x_0$ 时的无穷小量（无穷大量），且 $g(x)\ne0$，$x\in U_0(x_0,\delta_0)(\delta_0>0)$。

(1) 若 $\lim\limits_{x\to x_0}\dfrac{f(x)}{g(x)}=0$，则称 $f(x)$ 是比 $g(x)$ 更高阶的无穷小量（更低阶的无穷大量），记为 $f(x)=o(g(x))$（$x\to x_0$），并读做 $f(x)$ 等于小欧 $g(x)$；

(2) 若 $\lim\limits_{x\to x_0}\dfrac{f(x)}{g(x)}=l\ne0$，则称 $f(x)$ 与 $g(x)$ 是同阶无穷小量（同阶无穷大量）；

(3) 若 $\lim\limits_{x\to x_0}\dfrac{f(x)}{g(x)}=1$，则称 $f(x)$ 与 $g(x)$ 是等价无穷小量（等价无穷大量），记为 $f(x)\sim g(x)$（$x\to x_0$）；

(4) 若存在 $M>0$ 和 $\delta>0$，使得 $\forall x\in U_0(x_0,\delta)$，有 $|f(x)|\le M|g(x)|$ 成立，则记为 $f(x)=O(g(x))$（$x\to x_0$）（读做 $f(x)$ 等于大欧 $g(x)$）。

有的时候我们可以在上述定义中不假定 $g(x)$ 是一个无穷小（大）量. 如令 $g(x)\equiv1$，此时我们有以下记号
\[
  f(x)=O(1)\quad(x\to x_0).
\]
由 (4) 我们知道此记号表示 $f(x)$ 在 $x_0$ 的某个去心邻域上有界；而记号
\[
  f(x)=o(1)\quad(x\to x_0)
\]
表示 $\lim\limits_{x\to x_0}f(x)=0$，即 $f(x)$ 为 $x\to x_0$ 时的无穷小量。

在前面所给出的定义和记号中将“$x_0$”换成“$x_0+0$”，“$x_0-0$”，“$-\infty$”，“$+\infty$”或“$\infty$”，就得到相应的极限过程的定义和记号. 当然，此时对凡涉及“$x_0$ 的邻域”的地方都要做相应的修改。

特别地，我们可以对序列引进上述记号，如
\[
  \frac1{n^n}=o\left(\frac1n\right)\quad(n\to\infty).
\]

\noindent\textbf{例 3.4.1}\quad 函数 $y=x^n\sin\dfrac1x=O(x^n)$（$x\to0$）。

在学习无穷小量的比较时我们要注意以下的一个重要问题，即有时候不能简单地在两个无穷小量之间加上等号. 例如：从 $x^2=o(x)$（$x\to0$）我们可以推出 $x^2=O(x)$（$x\to0$）. 但我们不能简单地推出以下的等式：$o(x)=O(x)$（$x\to0$）. 读者很容易举出这种等式不成立的例子。

\noindent\textbf{例 3.4.2}\quad 设 $f(x)=(x-x_0)^m$，$g(x)=(x-x_0)^n$，则有
\[
  \lim_{x\to x_0}\frac{f(x)}{g(x)}
  =
  \begin{cases}
    0, & m>n,\\
    1, & m=n,\\
    \infty, & m<n.
  \end{cases}
\]
这说明：对于极限过程 $x\to x_0$，正整数 $k$ 越大时，无穷小量 $(x-x_0)^k$ 的阶就越高。

\noindent\textbf{例 3.4.3}\quad 设 $f(x)=\dfrac1{(x-x_0)^m}$，$g(x)=\dfrac1{(x-x_0)^n}$，则有
\[
  \lim_{x\to x_0}\frac{f(x)}{g(x)}
  =
  \begin{cases}
    0, & m<n,\\
    1, & m=n,\\
    \infty, & m>n.
  \end{cases}
\]
这说明：对于极限过程 $x\to x_0$，正整数 $k$ 越大时，无穷大量 $\dfrac1{(x-x_0)^k}$ 的阶就越高。

\noindent\textbf{例 3.4.4}\quad 设 $f(x)=x^\mu$（$\mu>0$），$g(x)=a^x$（$a>1$），证明
\[
  \lim_{x\to+\infty}\frac{f(x)}{g(x)}
  =
  \lim_{x\to+\infty}\frac{x^\mu}{a^x}=0.
\]

\noindent\textbf{证明}\quad 先考虑 $\mu=k$ 为正整数的情形. 令 $a=1+b$（$b>0$），则对任意的正整数 $n$，有
\[
  a^n=(1+b)^n
  =
  1+nb+\frac{n(n-1)}2b^2+\cdots+b^n
  >
  \frac{n(n-1)}2b^2,
\]
从而有
\[
  0<\frac n{a^n}<\frac2{(n-1)b^2},
\]
所以
\[
  \lim_{n\to\infty}\frac n{a^n}=0.
\]
由此立即可得
\[
  \lim_{n\to\infty}\frac{n^k}{a^n}
  =
  \lim_{n\to\infty}
  \left(\frac{n}{(\sqrt[k]{a})^n}\right)^k
  =
  0.
\]
这样，由于
\[
  0<\frac{x^k}{a^x}\le \frac{(1+[x])^k}{a^{[x]}},
\]
所以我们有
\[
  \lim_{x\to+\infty}\frac{x^k}{a^x}=0.
\]
对于一般的 $\mu>0$，我们取一个正整数 $k>\mu$，使得
\[
  0<\frac{x^\mu}{a^x}\le\frac{x^k}{a^x}\quad(x>1),
\]
从而有
\[
  \lim_{x\to+\infty}\frac{x^\mu}{a^x}=0.
\]
本例说明：对于极限过程 $x\to+\infty$，指数函数 $a^x$（$a>1$）是比任何幂函数 $x^\mu$ 都高阶的无穷大量。

\noindent\textbf{例 3.4.5}\quad 设 $f(x)=\log_a x$（$a>1$），$g(x)=x^\mu$（$\mu>0$），证明
\[
  \lim_{x\to+\infty}\frac{f(x)}{g(x)}
  =
  \lim_{x\to+\infty}\frac{\log_a x}{x^\mu}=0.
\]

\noindent\textbf{证明}\quad 令 $y=\log_a x$，则 $x=a^y$. 由此当 $x\to+\infty$ 时，$y\to+\infty$，有
\[
  \lim_{x\to+\infty}\frac{\log_a x}{x^\mu}
  =
  \lim_{y\to+\infty}\frac{y}{(a^\mu)^y}=0.
\]
本例说明：对于极限过程 $x\to+\infty$，对数函数 $\log_a x$（$a>1$）是比任何幂函数 $x^\mu$ 都低阶的无穷大量。

当 $x\to0$ 时，我们有以下重要的例子. 这些例子在今后的学习中会经常用到，读者应该熟练掌握它们。

\noindent\textbf{例 3.4.6}\quad 证明：当 $x\to0$ 时，下列关系式成立：
\[
  \text{(1) }\sin x\sim x;\qquad \sin x=x+o(x).
\]
\[
  \text{(2) }1-\cos x\sim\frac12x^2;\qquad
  \cos x=1-\frac12x^2+o(x^2).
\]
\[
  \text{(3) }\tan x\sim x;\qquad \tan x=x+o(x).
\]
\[
  \text{(4) }\arcsin x\sim x;\qquad \arcsin x=x+o(x).
\]
\[
  \text{(5) }\ln(1+x)\sim x;\qquad \ln(1+x)=x+o(x).
\]
\[
  \text{(6) }e^x-1\sim x;\qquad e^x=1+x+o(x).
\]
\[
  \text{(7) }(1+x)^\mu-1\sim\mu x;\qquad
  (1+x)^\mu=1+\mu x+o(x).
\]

\noindent\textbf{证明}\quad (1) $\lim\limits_{x\to0}\dfrac{\sin x}{x}=1$。

(2)
\[
  \lim_{x\to0}\frac{1-\cos x}{\frac12x^2}
  =
  \lim_{x\to0}\frac{2\sin^2\frac x2}{\frac12x^2}
  =
  \lim_{x\to0}\left(\frac{\sin\frac x2}{\frac x2}\right)^2
  =
  1.
\]

(3)
\[
  \lim_{x\to0}\frac{\tan x}{x}
  =
  \lim_{x\to0}\frac{\sin x}{x}\frac1{\cos x}
  =
  1.
\]

(4) 令 $t=\arcsin x$，则 $x=\sin t$，且当 $x\to0$ 时，$t\to0$. 于是，有
\[
  \lim_{x\to0}\frac{\arcsin x}{x}
  =
  \lim_{t\to0}\frac t{\sin t}
  =
  1.
\]

(5)
\[
  \lim_{x\to0}\frac{\ln(1+x)}x
  =
  \lim_{x\to0}\ln(1+x)^{1/x}
  =
  \ln e
  =
  1.
\]

(6) 令 $e^x-1=t$，则 $x=\ln(1+t)$，且当 $x\to0$ 时，$t\to0$. 因此
\[
  \lim_{x\to0}\frac{e^x-1}{x}
  =
  \lim_{t\to0}\frac t{\ln(1+t)}
  =
  1.
\]

(7)
\[
  \lim_{x\to0}\frac{(1+x)^\mu-1}{\mu x}
  =
  \lim_{x\to0}
  \left(\frac{e^{\mu\ln(1+x)}-1}{\mu\ln(1+x)}\right)
  \left(\frac{\ln(1+x)}x\right)
  =
  1.
\]

设 $f(x)$ 与 $g(x)$ 是当 $x\to x_0$ 时的等价无穷小量，则 $\dfrac{f(x)}{g(x)}-1=o(1)$（$x\to x_0$）. 由此推出 $f(x)=g(x)+o(g(x))$（$x\to x_0$），从而上述各结论的相应等式成立。

下面的定理说明，在求乘积或者商的极限时，可以将任何一个无穷小（大）的因式用它的等价因式来替换。

\noindent\textbf{定理 3.4.2}\quad 若无穷小（大）量 $h(x)\sim \ell(x)$（$x\to x_0$），则有
\[
  \text{(1) }\lim_{x\to x_0}h(x)f(x)=\lim_{x\to x_0}\ell(x)f(x);
\]
\[
  \text{(2) }\lim_{x\to x_0}\frac{h(x)f(x)}{g(x)}
  =
  \lim_{x\to x_0}\frac{\ell(x)f(x)}{g(x)};
\]
\[
  \text{(3) }\lim_{x\to x_0}\frac{f(x)}{h(x)g(x)}
  =
  \lim_{x\to x_0}\frac{f(x)}{\ell(x)g(x)}.
\]
这里，我们假定所有的函数都在点 $x_0$ 的某个去心邻域上有定义，作为分母的函数在这个去心领域上不为 $0$，并且假定各式右端的极限存在。

\noindent\textbf{证明}\quad 我们这里只给出结论 (1) 的证明，其余的请读者自己补出. 事实上，我们有
\[
  \lim_{x\to x_0}h(x)f(x)
  =
  \lim_{x\to x_0}\frac{h(x)}{\ell(x)}\ell(x)f(x)
  =
  \lim_{x\to x_0}\ell(x)f(x).
\]

\noindent\textbf{例 3.4.7}\quad 求当 $x\to0$ 时无穷小量 $\ln\cos 2x$ 的形如 $ax^\alpha$ 的等价无穷小量。

\noindent\textbf{解}\quad 由于
\[
  \ln\cos 2x
  =
  \ln(1-2\sin^2x)
  \sim
  -2\sin^2x
  \sim
  -2x^2,
\]
所以
\[
  \ln\cos 2x\sim -2x^2.
\]

\noindent\textbf{注}\quad 如果不是乘积和商的极限，则一般不能用等价无穷小（大）量来代替. 例如：
\[
  \lim_{x\to0}\frac{\tan x-\sin x}{x^3}
  =
  \frac12
  \ne
  \lim_{x\to0}\frac{x-x}{x^3}
  =
  0.
\]

\section*{习题三}

\noindent 1. 证明函数极限的定义与下述的 (1) 或 (2) 均等价：

(1) 对任意给定的正整数 $n$，存在正数 $\delta$，使当 $0<|x-x_0|<\delta$ 时，
\[
  |f(x)-A|<\frac1n;
\]

(2) 对任意给定的正数 $\varepsilon$，存在正整数 $n_0$，使当 $0<|x-x_0|<\dfrac1{n_0}$ 时，$|f(x)-A|<\varepsilon$。

\noindent 2. 用极限的定义证明下列极限：
\[
  \text{(1) }\lim_{x\to a}|x|=|a|;\qquad
  \text{(2) }\lim_{x\to3}x^3=27;\qquad
  \text{(3) }\lim_{x\to4}\sqrt{x}=2.
\]

\noindent 3. 叙述下列极限的定义：

\end{document}
